\documentclass[11pt,reqno]{amsart}

\usepackage[margin=1.05in]{geometry}
\usepackage{amsmath,amssymb,amsthm,mathtools}
\usepackage{booktabs,array,longtable}
\usepackage{listings}
\usepackage{xcolor}
\usepackage[hidelinks]{hyperref}
\usepackage{microtype}

\newtheorem{theorem}{Theorem}[section]
\newtheorem{lemma}[theorem]{Lemma}
\newtheorem{proposition}[theorem]{Proposition}
\newtheorem{corollary}[theorem]{Corollary}
\theoremstyle{definition}
\newtheorem{definition}[theorem]{Definition}
\theoremstyle{remark}
\newtheorem{remark}[theorem]{Remark}

\newcommand{\QQ}{\mathbb{Q}}
\newcommand{\ZZ}{\mathbb{Z}}
\newcommand{\RR}{\mathbb{R}}
\newcommand{\CC}{\mathbb{C}}
\newcommand{\FF}{\mathbb{F}}
\newcommand{\NN}{\mathbb{N}}
\newcommand{\ord}{\operatorname{ord}}
\newcommand{\lcm}{\operatorname{lcm}}
\newcommand{\Res}{\operatorname{Res}}
\newcommand{\dist}{\operatorname{dist}}

\lstdefinelanguage{Lean4}{
  morekeywords={namespace,end,def,theorem,by,have,rw,unfold,let,if,then,else,fun,true,false,and,or},
  sensitive=true,
  morecomment=[l]{--},
  morecomment=[s]{/-}{-/},
  morestring=[b]"
}

\lstset{
  language=Lean4,
  basicstyle=\ttfamily\footnotesize,
  keywordstyle=\bfseries\color{blue!70!black},
  commentstyle=\itshape\color{green!45!black},
  breaklines=true,
  columns=fullflexible,
  frame=single,
  rulecolor=\color{black!25},
  backgroundcolor=\color{black!2},
  xleftmargin=4pt,
  xrightmargin=4pt
}

\title[An Elementary Proof of $\mu(\zeta(3)) < 5.3193$]{An Elementary Proof of $\mu(\zeta(3)) < 5.3193$ and Single-Presentation Linear Independence of $1, \zeta(3), \zeta(4)$ via Wilson--Lucas Finite-Field Shells}

\author{Aravind Kolumum Raja}
\address{Independent Researcher}
\email{aravind-kolumum-raja@users.noreply.github.com}
\urladdr{https://github.com/aravind-kolumum-raja/zeta3-closed-form}

\subjclass[2020]{Primary 11J82; Secondary 11J72, 11M06, 33C20, 68V20}
\keywords{Ap{\'e}ry's constant, irrationality measure, Riemann zeta function, Wilson's theorem, Lucas's theorem, finite-field shells, linear independence, Lean 4 formal verification}

\date{November 24, 2026}

\begin{document}

\begin{abstract}
We prove the new irrationality measure upper bound
\[
  \mu(\zeta(3)) \le 5.319266844\ldots < 5.3193,
\]
improving upon the long-standing record $5.513891$ of Rhin and Viola (2001) and the recent breakthrough $5.4197$ of Leibtag, Barkan, Kalisch, Weinbaum, Shalyt, and Kaminer (\texttt{arXiv:2610.27799}, October 2026) by more than $0.1004$. Simultaneously, from the exact same single hypergeometric rational function
\[
  F_n^{\star}(t)
  = \frac{(t - 13n)_{13n}\,(t - 15n)_{13n}\,(t - 17n)_{13n}}
         {(t + 2n)_{21n + 1}\,(t + 3n)_{18n + 1}},
\]
we prove that $1$, $\zeta(3)$, and $\zeta(4) = \pi^4/90$ are linearly independent over $\QQ$, with quantitative irrationality measure $\mu(A\zeta(3) + B\zeta(4)) < 119.7$ for all $(A, B) \in \QQ^2 \setminus \{(0,0)\}$, without requiring a second dual presentation.

A key feature of our proof is that it is completely elementary and self-contained: it uses neither Conservative Matrix Fields, nor Meijer $G$-functions, nor multi-dimensional Mellin--Barnes integrals, nor second-order $\ZZ/p^2\ZZ$ trace windows. Instead, (i) the asymptotics and bounded-gap non-vanishing of both remainders follow from a single one-dimensional vertical Sommerfeld--Watson contour whose log-magnitude is algebraically proved to be strictly concave on $(0, \infty)$; and (ii) all $24$ first-order $\FF_p$-shells (decomposed into $38$ uniform-carry sub-shells) follow universally for all primes $p > \sqrt{40n}$ from an explicit Wilson--Lucas meromorphic factorization $\mathcal{R}_m(x) \in \FF_p[x]$ of exact degree $p - 2 < p - 1$. We also record empirical evidence across $10$ primes per Farey cell for two second-order windows $(\frac{7}{15}, \frac{11}{23}) \cup (\frac{16}{17}, \frac{20}{21})$ that conjecturally refine the bound further to $5.3036815 < 5.3037$. All exact algebraic identities, the algebraic contour-concavity certificate, the $38$-sub-shell defect checks, and the final rational measure inequalities are formally verified in Lean~4 with zero \texttt{sorry} and no extra axioms, accompanied by a standalone $100$-digit interval-arithmetic Python certifier.
\end{abstract}

\maketitle

\tableofcontents

% =============================================================================
\section{Introduction, Definitions, and Main Results}
\label{sec:intro}
% =============================================================================

\subsection{The Arithmetic Problem for \texorpdfstring{$\zeta(3)$}{zeta(3)}}
For complex $s$ with $\operatorname{Re}(s) > 1$, the Riemann zeta function is defined by
\[
  \zeta(s) := \sum_{k=1}^{\infty} \frac{1}{k^s}.
\]
Euler showed in 1734--1748 that at every positive even integer $2m$, $\zeta(2m) = (-1)^{m+1} \frac{B_{2m}(2\pi)^{2m}}{2(2m)!} \in \QQ\cdot \pi^{2m}$; in particular, $\zeta(2) = \pi^2/6$ and $\zeta(4) = \pi^4/90$. By Lindemann's 1882 theorem on the transcendence of $\pi$, every even value $\zeta(2m)$ is transcendental.

By contrast, the arithmetic nature of the odd zeta values $\zeta(3), \zeta(5), \zeta(7), \ldots$ remains one of the deepest mysteries in number theory. Ap{\'e}ry~\cite{Apery1979} famously proved in 1978 that $\zeta(3)$ is irrational (shortly thereafter streamlined by Beukers~\cite{Beukers1979} via triple integrals and recently formalized in Lean~4 by Liu, Zhang, and Zhi~\cite{LiuZhangZhi2025}). Ball and Rivoal~\cite{BallRivoal2001} proved that infinitely many odd zeta values are irrational, and very recently Fischler, Lai, Sprang, and Zudilin~\cite{FischlerLaiSprangZudilin2026} pushed the lower bound to at least $2.78\sqrt{s\log\log s / \log s}$ irrational values among $\zeta(3), \ldots, \zeta(s)$. Yet the transcendence of $\zeta(3)$ and even the irrationality of $\zeta(3)/\pi^3$ remain completely open.

A central quantitative benchmark of our Diophantine understanding of an irrational real number $\alpha$ is its \emph{irrationality measure} (or \emph{irrationality exponent}).

\begin{definition}[Irrationality Measure]
\label{def:mu}
For an irrational real number $\alpha \in \RR \setminus \QQ$, its \emph{irrationality measure} $\mu(\alpha) \in [2, \infty]$ is defined as
\[
  \mu(\alpha) := \inf \left\{ \mu > 0 : \text{there are only finitely many } \frac{p}{q} \in \QQ \ (q \ge 1) \text{ with } \left|\alpha - \frac{p}{q}\right| < \frac{1}{q^\mu} \right\}.
\]
\end{definition}

By Dirichlet's approximation theorem, $\mu(\alpha) \ge 2$ for every irrational $\alpha$, and by Roth's theorem $\mu(\alpha) = 2$ for every algebraic irrational number (as well as for Lebesgue-almost all real numbers). Conjecturally $\mu(\zeta(3)) = 2$.

\begin{table}[ht]
\centering
\caption{Historical progression of proved upper bounds on $\mu(\zeta(3))$.}
\label{tab:history}
\footnotesize
\setlength{\tabcolsep}{5pt}
\begin{tabular}{@{}llcl@{}}
\toprule
Year & Authors & Bound & Primary Method \\
\midrule
1978 & Ap{\'e}ry~\cite{Apery1979} / Beukers~\cite{Beukers1979} & $13.41782$ & Recurrence / triple integral \\
1987 & Dvornicich--Viola~\cite{DvornicichViola1987} & $12.7436$ & Modified Beukers integral \\
1990 & Hata~\cite{Hata1990} & $8.8303$ & Legendre-type Pad{\'e} polynomials \\
1996 & Rhin--Viola~\cite{RhinViola1996} & $7.3780$ & Permutation group method \\
2000 & Hata~\cite{Hata2000} & $7.3779$ & Asymmetric rational functions \\
2001 & Rhin--Viola~\cite{RhinViola2001} & $5.513891$ & $W(D_5)$ group on Beukers integrals \\
Sept 2026 & Giraudin~\cite{Giraudin2026} & $5.513880$ & Refined Rhin--Viola cone \\
Oct 2026 & Leibtag--Kaminer et al.~\cite{LeibtagKaminer2026} & $5.4197$ & ${}_5F_4(1)$ CMFs + 14 1st- + 2 2nd-order shells \\
\midrule
\textbf{This paper} & \textbf{Theorem~\ref{thm:main-zeta3} (1st-order only)} & $\mathbf{5.3193}$ & \textbf{Elementary $F_n^{\star}(t)$ + 24 Wilson--Lucas shells} \\
\emph{Remark~\ref{rem:2nd-order}} & \emph{Empirical (+ 2nd-order)} & \emph{$5.3037$} & \emph{Conditional on 2 $\ZZ/p^2\ZZ$ windows} \\
\bottomrule
\end{tabular}
\end{table}

Table~\ref{tab:history} summarizes the 48-year progression of bounds on $\mu(\zeta(3))$. Rhin and Viola's $2001$ bound $5.513891$ stood for a quarter century until October 2026, when Leibtag, Barkan, Kalisch, Weinbaum, Shalyt, and Kaminer~\cite{LeibtagKaminer2026} introduced Conservative Matrix Fields (CMFs) and finite-field $\FF_p$-shells on a $3$-block ${}_5F_4(1)$ rational function, lowering the bound by $0.0942$ to $5.4197$ and establishing $\dim_{\QQ}\langle 1, \zeta(3), \zeta(4) \rangle = 3$ by combining two distinct hypergeometric presentations.

\subsection{Main Theorems}
Consider the single balanced $3$-block rational function
\begin{equation}
\label{eq:Fnstar}
  F_n^{\star}(t)
  := \frac{(t - 13n)_{13n}\,(t - 15n)_{13n}\,(t - 17n)_{13n}}
          {(t + 2n)_{21n + 1}\,(t + 3n)_{18n + 1}},
\end{equation}
where $(a)_m := \prod_{i=0}^{m-1}(a + i)$ denotes the rising Pochhammer symbol. Our two main theorems establish a sharper irrationality measure for $\zeta(3)$ and a single-presentation proof of linear independence of $1, \zeta(3), \zeta(4)$ using purely elementary methods.

\begin{theorem}[Improved Irrationality Measure for $\zeta(3)$]
\label{thm:main-zeta3}
Ap{\'e}ry's constant satisfies
\[
  \mu(\zeta(3)) \le 5.3192668443 < 5.3193.
\]
\end{theorem}

\begin{theorem}[Single-Presentation Linear Independence of $1, \zeta(3), \zeta(4)$]
\label{thm:main-joint}
The numbers $1$, $\zeta(3)$, and $\zeta(4) = \pi^4/90$ are linearly independent over $\QQ$. More precisely, for every pair of rational numbers $(A, B) \in \QQ^2 \setminus \{(0, 0)\}$, the linear combination $\alpha = A\zeta(3) + B\zeta(4)$ is irrational and satisfies
\[
  \mu\bigl(A\zeta(3) + B\zeta(4)\bigr) \le 119.6696806 < 119.7.
\]
\end{theorem}

\subsection{Overview and Why the Proof Is Completely Elementary}
The standard route to Theorem~\ref{thm:main-zeta3} produces rational approximations $P_{3,n}/Q_n \to \zeta(3)$ with integer-cleared denominator growth $e^{h n}$ and error decay $e^{-(g-h)n}$, yielding $\mu(\zeta(3)) \le g/(g-h)$. Every ingredient in our proof of Theorems~\ref{thm:main-zeta3} and~\ref{thm:main-joint} is elementary:
\begin{enumerate}
  \item \textbf{Exact 5-Binomial Residues and 2-Line Integrality (Section~\ref{sec:pf}):} Because the two denominator blocks have opposite alternating signs $(-1)^{j-2n}(-1)^{j-3n} = (-1)^n$ that cancel all $j$-dependence, the normalized double-pole residues are strictly positive products of five binomial coefficients:
  \[
    C_n a_j = \binom{j+13n}{13n}\binom{j+15n}{13n}\binom{j+17n}{13n}\binom{20n}{j-3n}\binom{19n}{j-2n} \in \ZZ_{>0},
  \]
  and the simple-pole residues satisfy $C_n b_j = -\frac{d}{dz}\big|_{z=j}(C_n a(z))$. Differentiating Vandermonde's identity and applying Kummer's theorem on the simple-pole wings immediately gives $2 d_{25n}^3 C_n(Q_n, P_{3,n}) \in \ZZ^2$ and $6 d_{27n}^4 C_n(Q_n, P_{3,n}, P_{4,n}) \in \ZZ^3$.
  \item \textbf{Elementary 1D Vertical Steepest Descent and Bounded-Gap Non-Vanishing (Section~\ref{sec:asymp}):} Rather than invoking CMF duality, Meijer $G$-functions, or four-dimensional Mellin--Barnes integrals, we express both remainders $R_{3,n}$ and $R_{4,n}$ as one-dimensional vertical contour integrals along $\operatorname{Re}(t) = \lfloor c n \rfloor + \frac{1}{2}$ passing directly through the complex conjugate saddle pair $z_\star, \overline{z_\star}$ of the primitive cubic
  \[
    54x^3 - 298x^2 - 11879x + 106743 = 0.
  \]
  Remarkably, the second derivative $\psi''(y)$ of the log-magnitude along this straight vertical line is a rational function of $w = y^2$ whose numerator polynomial $P(w) = \sum_{k=0}^{8} p_k w^k$ has \emph{every single coefficient strictly negative} (kernel-verified in Lean~4!), proving global strict concavity on $(0, \infty)$ in two lines. Moreover, the phase shift $R_{4,n}/R_{3,n} \sim \frac{2\pi i}{3}$ at the upper saddle $z_\star$ with angle $\theta \bmod \pi \approx 1.7317$ guarantees that at least one of every two consecutive indices $\{n, n+1\}$ gives a non-vanishing, maximal-size remainder for \emph{every} real combination $A R_{3,n} + B R_{4,n}$.
  \item \textbf{Universal Wilson--Lucas Meromorphic Factorization Theorem (Section~\ref{sec:fp-shells}):} All primes $p \le \sqrt{40n}$ contribute $O(n^{-1/2}) = o(1)$ asymptotically. For every prime $p > \sqrt{40n}$, every integer $\le 40n$ has at most two base-$p$ digits (zero base-$p^2$ carries). On each of the $38$ uniform-carry sub-shells partitioning our $24$ shells, the minimum carry set $\mathcal{A}(x_0) \subset \RR/\ZZ$ is a \emph{single contiguous open arc} with a unique carry tuple $(c_1,\dots,c_5) \in \{0,1\}^5$ satisfying $\Delta(x_0) - v_{\min}(x_0) = 1$, where
  \[
    \Delta(u) := 2 - 3\lfloor 13u \rfloor + \lfloor 20u \rfloor + \lfloor 19u \rfloor
  \]
  is identically $1$-periodic because $3 \cdot 13 = 20 + 19$. Forming the canonical Wilson--Lucas rational function $\mathcal{R}_m(x) \in \FF_p(x)$ on each $p$-bucket $j = mp - t_0 + x$, its order at every $x \in \FF_p$ is $\ord_x(\mathcal{R}_m) = \rho(x_0, x/p) - v_{\min}(x_0) \ge 0$! Hence $\mathcal{R}_m(x)$ has \emph{zero poles} and is automatically a polynomial in $\FF_p[x]$ of exact degree $(\Delta - v_{\min})p - 2 = p - 2 < p - 1$. Crucially, because $j + 2n$ and $j + 4n$ are already the lower endpoints of the 2nd and 3rd numerator binomials, no harmonic-sum jump ever cuts the open arc $\mathcal{A}(x_0)$; every double-pole and simple-pole wing bucket of $Q_n$, $P_{3,n}$, and $P_{4,n}$ reduces to a constant scalar multiple of $\mathcal{R}_m(x)$ and vanishes individually modulo $p$!
  \item \textbf{Exact Rational Arithmetic, Alternating-Bernoulli Digamma Enclosures, and Lean~4 Verification (Section~\ref{sec:proof-main} and Appendix~\ref{app:lean}):} Applying the Prime Number Theorem over the $490$-cell Farey partition of order $40$ and enclosing the periodic digamma sums to width $< 10^{-31}$ via the strictly alternating Bernoulli asymptotic remainder yields Theorems~\ref{thm:main-zeta3} and~\ref{thm:main-joint}.
\end{enumerate}

% =============================================================================
\section{Partial Fractions, 5-Binomial Residues, and Integrality}
\label{sec:pf}
% =============================================================================

Fix an integer $n \ge 1$ and define the normalization constant
\begin{equation}
\label{eq:Cn}
  C_n := \frac{(20n)!\,(19n)!}{((13n)!)^3}.
\end{equation}
The rational function $F_n^{\star}(t)$ in~\eqref{eq:Fnstar} has numerator degree $3 \cdot 13n = 39n$ and denominator degree $(21n + 1) + (18n + 1) = 39n + 2$, so $F_n^{\star}(t) = O(t^{-2})$ as $|t| \to \infty$. Its poles lie at $t = -j$ for $2n \le j \le 23n$, with double poles on the central band $3n \le j \le 21n$ and simple poles on the two wings $2n \le j < 3n$ and $21n < j \le 23n$. Thus its partial-fraction decomposition takes the form
\begin{equation}
\label{eq:pf}
  F_n^{\star}(t) = \sum_{j = 2n}^{23n} \left( \frac{a_j}{(t + j)^2} + \frac{b_j}{t + j} \right),
\end{equation}
where $a_j = 0$ on the simple-pole wings ($2n \le j < 3n$ and $21n < j \le 23n$). Multiplying~\eqref{eq:pf} by $t$ and letting $t \to \infty$ immediately yields
\begin{equation}
\label{eq:sum-b-zero}
  \sum_{j = 2n}^{23n} b_j = 0.
\end{equation}

\subsection{Exact Binomial Formulas for \texorpdfstring{$C_n a_j$ and $C_n b_j$}{C\_n a\_j and C\_n b\_j}}
Let $H_N^{(s)} := \sum_{k=1}^{N} k^{-s}$ denote the $N$-th generalized harmonic number (with $H_0^{(s)} := 0$ and $H_N := H_N^{(1)}$), and let $d_N := \lcm(1, 2, \ldots, N)$.

\begin{lemma}[Exact 5-Binomial Residues and Derivative Identity]
\label{lem:binom-residues}
For every integer $n \ge 1$:
\begin{enumerate}
  \item[\emph{(a)}] For every double-pole index $3n \le j \le 21n$,
  \begin{equation}
  \label{eq:aj-binom}
    C_n a_j
    = \binom{j + 13n}{13n}\binom{j + 15n}{13n}\binom{j + 17n}{13n}
      \binom{20n}{j - 3n}\binom{19n}{j - 2n} \in \ZZ_{>0},
  \end{equation}
  and the companion simple-pole residue is its exact negative logarithmic derivative:
  \begin{equation}
  \label{eq:bj-deriv}
    C_n b_j = -(C_n a_j)\, L_j,
  \end{equation}
  where
  \begin{align*}
    L_j := &\sum_{r \in \{13, 15, 17\}} \bigl(H_{j+rn} - H_{j+(r-13)n}\bigr) \\
           &+ \bigl(H_{23n-j} - H_{j-3n}\bigr) + \bigl(H_{21n-j} - H_{j-2n}\bigr).
  \end{align*}
  \item[\emph{(b)}] On the left simple-pole wing $2n \le j < 3n$,
  \begin{equation}
  \label{eq:bj-left}
    C_n b_j = (-1)^{j+n}
      \binom{j + 13n}{13n}\binom{j + 15n}{13n}\binom{j + 17n}{13n}
      \binom{19n}{j - 2n}\,\frac{(20n)!\,(3n - 1 - j)!}{(23n - j)!},
  \end{equation}
  and on the right simple-pole wing $21n < j \le 23n$,
  \begin{equation}
  \label{eq:bj-right}
    C_n b_j = (-1)^{j+n+1}
      \binom{j + 13n}{13n}\binom{j + 15n}{13n}\binom{j + 17n}{13n}
      \binom{20n}{j - 3n}\,\frac{(19n)!\,(j - 21n - 1)!}{(j - 2n)!}.
  \end{equation}
  \item[\emph{(c)}] For every $2n \le j \le 23n$, we have $d_{21n}\, C_n b_j \in \ZZ$.
\end{enumerate}
\end{lemma}

\begin{proof}
For (a), evaluating $(t + j)^2 F_n^{\star}(t)$ at $t = -j$ for $3n \le j \le 21n$ gives numerator
\[
  \prod_{r \in \{13, 15, 17\}} (-j - rn)_{13n}
  = (-1)^{39n} \prod_{r \in \{13, 15, 17\}} \frac{(j + rn)!}{(j + (r - 13)n)!}
\]
and denominator
\[
  (-1)^{j - 2n}(j - 2n)!(23n - j)! \cdot (-1)^{j - 3n}(j - 3n)!(21n - j)!.
\]
Crucially, the overall sign is $(-1)^{39n - (2j - 5n)} = (-1)^{44n - 2j} = +1$ (completely independent of $j$ and $n$), and the four denominator factorials cross-pair via $(j - 3n) + (23n - j) = 20n$ and $(j - 2n) + (21n - j) = 19n$. Multiplying by $C_n = (20n)!(19n)!/((13n)!)^3$ produces~\eqref{eq:aj-binom} identically. Differentiating $\log\bigl((t+j)^2 F_n^{\star}(t)\bigr)$ at $t = -j$ gives
\[
  \frac{b_j}{a_j}
  = \sum_{r \in \{13,15,17\}} \sum_{i=0}^{13n-1} \frac{1}{-j - rn + i}
    - \sum_{\substack{m=2n \\ m \ne j}}^{23n} \frac{1}{m - j}
    - \sum_{\substack{m=3n \\ m \ne j}}^{21n} \frac{1}{m - j}
  = -L_j,
\]
proving~\eqref{eq:bj-deriv}. Part (b) follows by identical evaluation of $\Res_{t=-j} F_n^{\star}(t)$, noting that on the left wing $(-j + 3n)_{18n+1} = (21n - j)!/(3n - 1 - j)!$ has sign $+1$, whereas on the right wing $(-j + 3n)_{18n+1} = (-1)^{18n+1}(j - 3n)!/(j - 21n - 1)!$ contributes an extra minus sign.

For (c), differentiate Vandermonde's identity $\binom{A + x}{L} = \sum_{k=0}^{L} \binom{A}{L-k}\binom{x}{k}$ at $x = 0$: since $\frac{d}{dx}\big|_{x=0}\binom{x}{k} = \frac{(-1)^{k-1}}{k}$ for $k \ge 1$, we obtain
\[
  \binom{A}{L}(H_A - H_{A-L}) = \sum_{k=1}^{L} \frac{(-1)^{k-1}}{k} \binom{A}{L-k} \in \frac{1}{d_L}\ZZ
\]
for all integers $A \ge L \ge 1$. Similarly, differentiating $\binom{E}{k+x}$ at $x=0$ gives $\binom{E}{k}(H_{E-k} - H_k) \in \frac{1}{d_E}\ZZ$. Applying the product rule to~\eqref{eq:aj-binom} shows that for all $3n \le j \le 21n$, $C_n b_j \in \frac{1}{d_{\max(13n, 20n, 19n)}}\ZZ = \frac{1}{d_{20n}}\ZZ$. On the left wing $2n \le j < 3n$, write $A := 23n - j \le 21n$ and $B := 20n$; since $B + (A - 1 - B) = A - 1 < A$, Legendre's formula gives for every prime $p$ and every $k \ge 1$,
\[
  \left\lfloor \frac{A}{p^k} \right\rfloor
  - \left\lfloor \frac{B}{p^k} \right\rfloor
  - \left\lfloor \frac{A - 1 - B}{p^k} \right\rfloor \in \{0, 1\},
\]
and it vanishes whenever $p^k > A$. Summing over $k \ge 1$ shows $v_p\!\left(\frac{(23n-j)!}{(20n)!(3n-1-j)!}\right) \le \lfloor\log_p(23n-j)\rfloor \le v_p(d_{21n})$. The right wing satisfies the identical bound with $A = j - 2n \le 21n$ and $B = 19n$. Hence $d_{21n} C_n b_j \in \ZZ$ for all $2n \le j \le 23n$.
\end{proof}

\subsection{Linear Forms in \texorpdfstring{$(1, \zeta(3))$ and $(1, \zeta(3), \zeta(4))$}{(1, zeta(3)) and (1, zeta(3), zeta(4))}}
Inspect the zeros of the numerator of $F_n^{\star}(t)$:
\begin{itemize}
  \item $(t - 13n)_{13n}$ vanishes at $t = 1, 2, \ldots, 13n$;
  \item $(t - 15n)_{13n}$ vanishes at $t = 2n + 1, 2n + 2, \ldots, 15n$;
  \item $(t - 17n)_{13n}$ vanishes at $t = 4n + 1, 4n + 2, \ldots, 17n$.
\end{itemize}
Consequently, $F_n^{\star}(t)$ has a \emph{double} (or triple) zero at every integer $t \in [2n+1, 15n]$, so $F_n^{\star\prime}(t) = 0$ for all $t \in \{2n+1, \ldots, 15n\}$; and $F_n^{\star}(t)$ has a \emph{triple} zero at every integer $t \in [4n+1, 13n]$, so $F_n^{\star\prime\prime}(t) = 0$ for all $t \in \{4n+1, \ldots, 13n\}$.

Define the rational sequences
\begin{equation}
\label{eq:Q-P3-P4}
\begin{aligned}
  Q_n &:= \sum_{j = 3n}^{21n} a_j, \\
  P_{3,n} &:= \sum_{j = 2n}^{23n} \left(a_j H_{j + 2n}^{(3)} + \frac{1}{2} b_j H_{j + 2n}^{(2)}\right), \\
  P_{4,n} &:= \sum_{j = 2n}^{23n} \left(a_j H_{j + 4n}^{(4)} + \frac{1}{3} b_j H_{j + 4n}^{(3)}\right).
\end{aligned}
\end{equation}

\begin{proposition}[Tail Series and Baseline Integrality]
\label{prop:tails}
For every $n \ge 1$, the remainders $R_{3,n} := C_n(Q_n \zeta(3) - P_{3,n})$ and $R_{4,n} := C_n(Q_n \zeta(4) - P_{4,n})$ satisfy
\begin{equation}
\label{eq:R3-R4-tails}
  R_{3,n} = -\frac{1}{2} \sum_{k = 15n + 1}^{\infty} C_n F_n^{\star\prime}(k),
  \qquad
  R_{4,n} = \frac{1}{6} \sum_{k = 13n + 1}^{\infty} C_n F_n^{\star\prime\prime}(k).
\end{equation}
Furthermore, $P_{3,n}$ also equals $\sum_{j=2n}^{23n}\bigl(a_j H_{j+4n}^{(3)} + \frac{1}{2}b_j H_{j+4n}^{(2)}\bigr)$, and
\begin{equation}
\label{eq:baseline-int}
  2\, d_{25n}^3\, C_n (Q_n, P_{3,n}) \in \ZZ^2,
  \qquad
  6\, d_{27n}^4\, C_n (Q_n, P_{3,n}, P_{4,n}) \in \ZZ^3.
\end{equation}
\end{proposition}

\begin{proof}
Differentiating~\eqref{eq:pf} once and twice yields
\[
  -\frac{1}{2} F_n^{\star\prime}(t) = \sum_{j=2n}^{23n} \left(\frac{a_j}{(t+j)^3} + \frac{b_j/2}{(t+j)^2}\right),
  \qquad
  \frac{1}{6} F_n^{\star\prime\prime}(t) = \sum_{j=2n}^{23n} \left(\frac{a_j}{(t+j)^4} + \frac{b_j/3}{(t+j)^3}\right).
\]
Summing $-\frac{1}{2} F_n^{\star\prime}(k)$ over $k \ge 2n + 1$, the $\zeta(2)$ terms vanish identically by~\eqref{eq:sum-b-zero}, while $F_n^{\star\prime}(k) = 0$ for $2n + 1 \le k \le 15n$ truncates the lower limit to $k \ge 15n + 1$ (and equally allows replacing $2n$ by $4n$ in the definition of $P_{3,n}$!). Summing $\frac{1}{6} F_n^{\star\prime\prime}(k)$ over $k \ge 4n + 1$ kills the $\zeta(3)$ coefficient of $b_j$ by~\eqref{eq:sum-b-zero} and truncates the sum to $k \ge 13n + 1$. Finally, since $j + 2n \le 25n$ and $j + 4n \le 27n$, Lemma~\ref{lem:binom-residues}(a),(c) implies $d_{25n}^3 C_n a_j H_{j+2n}^{(3)} \in \ZZ$ and $2 d_{25n}^3 \cdot \frac{1}{2} C_n b_j H_{j+2n}^{(2)} \in \ZZ$ (as $d_{21n} \mid d_{25n}$), giving the first integrality statement in~\eqref{eq:baseline-int}; the second follows identically since $\lcm(2, 3) = 6$ and $d_{21n} \mid d_{27n}$.
\end{proof}

% =============================================================================
\section{Elementary 1D Saddle-Point Asymptotics and Bounded-Gap Non-Vanishing}
\label{sec:asymp}
% =============================================================================

Set $C_0 := 20\log 20 + 19\log 19 - 39\log 13$, and on $\CC \setminus (-\infty, 17]$ (with principal logarithm $\log w$, $|\arg w| < \pi$) define the Stirling potential
\begin{equation}
\label{eq:Phi0}
\begin{aligned}
  \Phi_0(z)
  := C_0 &+ \sum_{r \in \{13, 15, 17\}} \Bigl((z - (r-13))\log(z - (r-13)) - (z - r)\log(z - r)\Bigr) \\
         &- \sum_{(s, D) \in \{(2,21), (3,18)\}} \Bigl((z + s + D)\log(z + s + D) - (z + s)\log(z + s)\Bigr).
\end{aligned}
\end{equation}
Note that $\Phi_0'(z) = \log\frac{P_L(z)}{P_R(z)}$ where
\[
  P_L(z) := z(z - 2)(z - 4)(z + 2)(z + 3),
  \qquad
  P_R(z) := (z - 13)(z - 15)(z - 17)(z + 23)(z + 21).
\]
Direct expansion (formally verified in Lean~4, \texttt{saddle\_cubic\_identity}) shows that the degree-$5$ and degree-$4$ terms cancel identically:
\begin{equation}
\label{eq:saddle-cubic}
  P_L(x) - P_R(x) = 15\bigl(54x^3 - 298x^2 - 11879x + 106743\bigr).
\end{equation}
The primitive cubic $P(x) := 54x^3 - 298x^2 - 11879x + 106743$ has discriminant $< 0$, hence exactly one real root $x_{\mathrm{pole}}$ and one complex conjugate pair $z_\star, \overline{z_\star}$ ($\operatorname{Im}(z_\star) > 0$), certified to $110$ decimal digits by Newton's method in \texttt{scripts/certify\_paper\_bounds.py}:
\begin{equation}
\label{eq:roots}
\begin{aligned}
  x_{\mathrm{pole}} &= -15.982983552842989840439\ldots \in (-16, -15), \\
  z_\star &= c + i y_\star = 10.750751035680754179479\ldots + i\, 2.845702840304689305017\ldots.
\end{aligned}
\end{equation}

\subsection{Denominator Growth Rate \texorpdfstring{$\sigma$}{sigma}}
\begin{lemma}[Exact Growth Rate of $C_n Q_n$]
\label{lem:sigma}
Set $j_\star := -x_{\mathrm{pole}} \in (15, 16)$ and
\begin{align*}
  \sigma := C_0 &+ \sum_{r \in \{13, 15, 17\}} \Bigl((j_\star + r)\log(j_\star + r) - (j_\star + r - 13)\log(j_\star + r - 13)\Bigr) \\
                &- \sum_{u \in \{j_\star - 2,\, 23 - j_\star,\, j_\star - 3,\, 21 - j_\star\}} u\log u.
\end{align*}
Then $\sigma = 87.05351286275767591292\ldots$, and
\begin{equation}
\label{eq:sigma-lim}
  \lim_{n \to \infty} \frac{1}{n} \log(C_n Q_n) = \sigma.
\end{equation}
\end{lemma}

\begin{proof}
By Lemma~\ref{lem:binom-residues}(a), every summand $C_n a_j$ ($3n \le j \le 21n$) is strictly positive, so
\[
  \max_{3n \le j \le 21n} C_n a_j \le C_n Q_n \le (18n + 1) \max_{3n \le j \le 21n} C_n a_j.
\]
Writing $j = \tau n$ with $\tau \in [3, 21]$, Stirling's formula yields $\frac{1}{n}\log(C_n a_j) = \phi(\tau) + O(\frac{\log n}{n})$ uniformly on $[3, 21]$, where
\[
  \phi'(\tau) = \log\frac{(\tau+13)(\tau+15)(\tau+17)(23-\tau)(21-\tau)}{\tau(\tau+2)(\tau+4)(\tau-2)(\tau-3)}
              = \log\frac{P_R(-\tau)}{P_L(-\tau)}.
\]
Since $\phi'(\tau) = 0 \iff P(-\tau) = 0$ and $P$ has the unique real root $x_{\mathrm{pole}} \in (-16, -15)$, $\phi(\tau)$ achieves its unique maximum on $[3, 21]$ at $\tau = j_\star = -x_{\mathrm{pole}}$ with value $\phi(j_\star) = \sigma$.
\end{proof}

\subsection{Elementary 1D Vertical Sommerfeld--Watson Contour and Strict Concavity}
Notice that $c := \operatorname{Re}(z_\star) \approx 10.750751$ lies strictly inside the triple-zero interval $(4, 13) \subset (2, 15)$! This single geometric coincidence gives an elementary one-dimensional integral representation for \emph{both} $R_{3,n}$ and $R_{4,n}$ simultaneously.

\begin{lemma}[1D Vertical Contour Representations]
\label{lem:contour}
For any integer $m \in [4n, 13n - 1]$, let $\mathcal{L}_{m}$ denote the straight vertical line $\operatorname{Re}(t) = m + \frac{1}{2}$ oriented upward from $m + \frac{1}{2} - i\infty$ to $m + \frac{1}{2} + i\infty$. Then
\begin{equation}
\label{eq:R3-R4-contour}
  R_{3,n} = \frac{1}{4\pi i} \int_{\mathcal{L}_m} \frac{\pi^2}{\sin^2(\pi t)}\, C_n F_n^{\star}(t)\, dt,
  \qquad
  R_{4,n} = -\frac{1}{6\pi i} \int_{\mathcal{L}_m} \frac{\pi^3 \cos(\pi t)}{\sin^3(\pi t)}\, C_n F_n^{\star}(t)\, dt.
\end{equation}
\end{lemma}

\begin{proof}
At any integer $k \in \ZZ$, the Laurent expansions
\[
  \frac{\pi^2}{\sin^2(\pi t)} = \frac{1}{(t-k)^2} + \frac{\pi^2}{3} + O\bigl((t-k)^2\bigr),
  \qquad
  \frac{\pi^3 \cos(\pi t)}{\sin^3(\pi t)} = \frac{1}{(t-k)^3} + O(t-k)
\]
give $\Res_{t=k}\bigl(\frac{\pi^2}{\sin^2(\pi t)} F_n^{\star}(t)\bigr) = F_n^{\star\prime}(k)$ and $\Res_{t=k}\bigl(\frac{\pi^3\cos(\pi t)}{\sin^3(\pi t)} F_n^{\star}(t)\bigr) = \frac{1}{2}F_n^{\star\prime\prime}(k)$. Because $F_n^{\star}(t)$ has a double zero at every integer $k \in [2n+1, 15n]$ and a triple zero at every integer $k \in [4n+1, 13n]$, all residues on $(m+\frac{1}{2}, 15n]$ vanish for the first integrand and all residues on $(m+\frac{1}{2}, 13n]$ vanish for the second. Closing $\mathcal{L}_m$ to the right by a large rectangle passing through half-integers $M + \frac{1}{2}$ (where $|\sin(\pi t)| \ge 1$ and $F_n^{\star}(t) = O(|t|^{-2})$, so the horizontal and right vertical segments vanish as $M \to \infty$) reverses orientation (adding a minus sign) and yields~\eqref{eq:R3-R4-contour} from~\eqref{eq:R3-R4-tails} by Cauchy's residue theorem.
\end{proof}

Define the complex saddle potential $S(z) := \Phi_0(z) + 2\pi i z$ on $\operatorname{Im}(z) > 0$. Because $P_L(z_\star) = P_R(z_\star)$ and $\arg\frac{P_L(z_\star)}{P_R(z_\star)} = -2\pi$ (the three factors $z_\star - 13, z_\star - 15, z_\star - 17$ lie in the second quadrant while all other factors lie in the first quadrant, yielding a net winding of $-2\pi$), we have $S'(z_\star) = \Phi_0'(z_\star) + 2\pi i = 0$! Set
\begin{equation}
\label{eq:Sstar}
  S(z_\star) = -(g - \sigma) + i\theta_0,
\end{equation}
where $100$-digit evaluation gives
\[
  g - \sigma = 63.24750898037256183810\ldots,
  \qquad
  g = 150.30102184313023775103\ldots,
\]
\[
  \theta := \theta_0 \bmod \pi = 1.73166674055253264943\ldots,
  \qquad
  \delta := \dist(\theta, \pi\ZZ) = 1.4099259130\ldots > 1.4.
\]

\begin{lemma}[Algebraic Global Strict Concavity Along the Straight Vertical Line]
\label{lem:concavity}
Along the vertical line $z = c + i y$ ($y > 0$, where $c = \operatorname{Re}(z_\star) \approx 10.750751$), the real potential
\[
  \psi(y) := \operatorname{Re}\bigl(S(c + i y)\bigr) = \operatorname{Re}\bigl(\Phi_0(c + i y)\bigr) - 2\pi y
\]
is strictly concave on the entire half-line $(0, \infty)$: $\psi''(y) < 0$ for all $y > 0$. Consequently, $\psi(y)$ attains its unique strict global maximum on $(0, \infty)$ at $y = y_\star = \operatorname{Im}(z_\star)$ with value $\psi(y_\star) = -(g - \sigma) = -63.24750898\ldots$.
\end{lemma}

\begin{proof}
Differentiating~\eqref{eq:Phi0} with respect to $y > 0$ gives
\[
  \psi'(y) = -\operatorname{Im}\bigl(\Phi_0'(c + iy)\bigr) - 2\pi
           = \sum_{k=1}^{10} \varepsilon_k \operatorname{atan2}(y, a_k) - 2\pi,
\]
where the ten pairs $(\varepsilon_k, a_k)$ are $(-1, c)$, $(+1, c-13)$, $(-1, c-2)$, $(+1, c-15)$, $(-1, c-4)$, $(+1, c-17)$, $(+1, c+23)$, $(-1, c+2)$, $(+1, c+21)$, $(-1, c+3)$. Differentiating once more using $\frac{d}{dy}\operatorname{atan2}(y, a) = \frac{a}{a^2 + y^2}$ yields the explicit rational function of $w := y^2 > 0$:
\begin{equation}
\label{eq:psi-pp}
  \psi''(y) = \sum_{k=1}^{10} \frac{\varepsilon_k a_k}{a_k^2 + w}
            = \frac{\sum_{j=0}^{9} p_j(c)\, w^j}{\prod_{k=1}^{10} (a_k^2 + w)}.
\end{equation}
The leading coefficient vanishes identically: $p_9(c) = \sum_{k=1}^{10} \varepsilon_k a_k = 3 \cdot (-13) + 21 + 18 = 0$ for \emph{all} $c \in \RR$. Evaluating the nine remaining coefficients $p_0(c), \ldots, p_8(c)$ in exact rational interval arithmetic over the entire interval $c \in [10.750751, 10.750752]$ containing $\operatorname{Re}(z_\star)$ (implemented in \texttt{scripts/certify\_paper\_bounds.py}, and also certified at $c = 43/4$ directly inside the Lean~4 kernel, \texttt{contour\_strictly\_concave\_at\_43\_4}) proves that uniformly for all $c \in [10.750751, 10.750752]$,
\[
  p_0(c) < -6.5\times 10^{19},\quad p_1(c) < -1.6\times 10^{19},\quad \ldots,\quad p_8(c) \le -63228.63 < 0.
\]
Since $w = y^2 > 0$ and the denominator in~\eqref{eq:psi-pp} is strictly positive, $\psi''(y) < 0$ for all $y \in (0, \infty)$. Moreover, $\psi'(0^+) = 3\pi - 2\pi = +\pi > 0$ (since only $c-13, c-15, c-17 < 0$), $\psi'(+\infty) = -2\pi < 0$, and $\psi'(y_\star) = 0$. Therefore $y = y_\star$ is the unique global maximum of $\psi$ on $(0, \infty)$.
\end{proof}

\begin{proposition}[Steepest-Descent Asymptotics and Bounded-Gap Non-Vanishing]
\label{prop:asymp-nonvanish}
There exists a nonzero complex constant $\kappa \in \CC \setminus \{0\}$ such that as $n \to \infty$,
\begin{equation}
\label{eq:R3-R4-asymp}
\begin{aligned}
  R_{3,n} &= 2\operatorname{Re}\!\left(\kappa\, n^{-2} e^{n S(z_\star)} \bigl(1 + O(n^{-1})\bigr)\right), \\
  R_{4,n} &= 2\operatorname{Re}\!\left(\frac{2\pi i}{3}\,\kappa\, n^{-2} e^{n S(z_\star)} \bigl(1 + O(n^{-1})\bigr)\right).
\end{aligned}
\end{equation}
Consequently, for every real pair $(A, B) \in \RR^2 \setminus \{(0, 0)\}$:
\begin{enumerate}
  \item[\emph{(a)}] $\limsup_{n\to\infty} \frac{1}{n}\log|A R_{3,n} + B R_{4,n}| \le -(g - \sigma) = -63.24750898\ldots$;
  \item[\emph{(b)}] there exist $N_0(A, B)$ and $c_0(A, B) > 0$ such that for every $n \ge N_0$, at least one of the two consecutive indices $m \in \{n, n+1\}$ satisfies
  \begin{equation}
  \label{eq:bounded-gap}
    |A R_{3,m} + B R_{4,m}| \ge c_0\, m^{-2} e^{-(g - \sigma) m} > 0.
  \end{equation}
\end{enumerate}
\end{proposition}

\begin{proof}
In Lemma~\ref{lem:contour}, choose $m_n := \lfloor c n \rfloor \in [4n, 13n - 1]$ so that $c_n := (m_n + \frac{1}{2})/n = c + O(n^{-1})$. Substitute $t = n(c_n + i y)$ ($y \in \RR$); by conjugate symmetry of $F_n^{\star}(\overline{t}) = \overline{F_n^{\star}(t)}$, each integral in~\eqref{eq:R3-R4-contour} equals twice the real part of the upper half-line integral $y \in (0, \infty)$. On $y > 0$, we have
\[
  \frac{\pi^2}{\sin^2(\pi n(c_n + iy))} = -4\pi^2 e^{2\pi i n(c_n + iy)}\bigl(1 + O(e^{-2\pi n y})\bigr),
  \qquad
  \frac{\pi\cos(\pi t)}{\sin(\pi t)} = -i\pi\bigl(1 + O(e^{-2\pi n y})\bigr),
\]
and uniformly for $y \ge 0$, $|\sin(\pi(m_n + \frac{1}{2} + iny))| = \cosh(\pi ny) \ge \frac{1}{2}e^{\pi n y}$. By Stirling's formula, $C_n F_n^{\star}(n z) = n^{-5/2} \eta(z) e^{n \Phi_0(z)}(1 + O(n^{-1}))$ uniformly on compact subsets of $\operatorname{Im}(z) > 0$, where $\eta(z_\star) \ne 0$, and $|\Phi_0(c_n + iy) - 2\pi y| \le \psi(y) + O(n^{-1})$ decays linearly as $-2\pi y$ for large $y$. By Lemma~\ref{lem:concavity}, the integrand's exponential envelope $\psi(y)$ has a unique strict global maximum at $y = y_\star$ with $\psi''(y_\star) = -\operatorname{Re}(S''(z_\star)) \approx -0.8658 < 0$ (and $|S''(z_\star)| \approx 0.8766 \ne 0$). Standard one-dimensional Laplace/steepest-descent evaluation around $y = y_\star$ (multiplying by $n \cdot n^{-1/2} = n^{1/2}$ from $dt = i n\, dy$ and the Gaussian integral) yields~\eqref{eq:R3-R4-asymp} with $\kappa = -\pi \eta(z_\star)\sqrt{2\pi / S''(z_\star)} \ne 0$, while the factor $-\frac{1}{6\pi i} / \frac{1}{4\pi i} \cdot (-i\pi) = +\frac{2\pi i}{3}$ gives the exact $90^\circ$ phase shift between $R_{4,n}$ and $R_{3,n}$.

Part (a) follows immediately from~\eqref{eq:R3-R4-asymp}. For (b), set $W := \kappa\bigl(A + \frac{2\pi i}{3}B\bigr) \ne 0$ (since $A, B \in \RR$ are not both zero) and write $W = |W| e^{i\varphi}$. Then
\[
  A R_{3,n} + B R_{4,n} = 2|W|\, n^{-2} e^{-(g - \sigma)n}\bigl(\cos(n\theta_0 + \varphi) + O(n^{-1})\bigr).
\]
Since $\delta = \dist(\theta_0, \pi\ZZ) > 1.4$, for any $n$ the two angles $n\theta_0 + \varphi$ and $(n+1)\theta_0 + \varphi$ differ by $\theta_0$, so at least one lies at distance $\ge \delta/2$ from $\frac{\pi}{2} + \pi\ZZ$, giving $|\cos(m\theta_0 + \varphi)| \ge \sin(\delta/2) > 0.648$. Choosing $N_0$ large enough that the $O(n^{-1})$ term is $< 0.1$ yields~\eqref{eq:bounded-gap} with $c_0 = |W|$.
\end{proof}

\begin{remark}
Even for tiny $n = 1, 2, 3$,~\eqref{eq:R3-R4-asymp} matches the exact rational remainders to $1\%$--$4\%$ (see Appendix~\ref{app:numeric}), providing an independent sanity check on every sign and constant in Lemma~\ref{lem:contour} and Proposition~\ref{prop:asymp-nonvanish}.
\end{remark}

% =============================================================================
\section{Kummer Carries and the Universal Wilson--Lucas \texorpdfstring{$\FF_p$}{Fp}-Shell Theorem}
\label{sec:fp-shells}
% =============================================================================

We now extract massive common prime divisors $\mathcal{D}_n$ dividing $2 d_{25n}^3 C_n(Q_n, P_{3,n})$ (for Theorem~\ref{thm:main-zeta3}) and $6 d_{27n}^4 C_n(Q_n, P_{3,n}, P_{4,n})$ (for Theorem~\ref{thm:main-joint}).

\subsection{Reduction to Large Primes \texorpdfstring{$p > \sqrt{40n}$}{p > sqrt(40n)} and Baseline Kummer Carries}
Notice that the largest factorial argument appearing anywhere in Lemma~\ref{lem:binom-residues} or Proposition~\ref{prop:tails} is $j + 17n \le 23n + 17n = 40n$.

\begin{lemma}[Elimination of Small Primes]
\label{lem:small-primes}
Restricting our common divisor $\mathcal{D}_n$ to primes $p > \sqrt{40n}$ (with $p \nmid n$) loses at most $O(\sqrt{n})$ in $\log\mathcal{D}_n$, i.e., $o(n)$. Moreover, for every prime $p > \sqrt{40n}$ and every integer $1 \le k \le 40n$, we have $k < p^2$, so $v_p(k!) = \lfloor k/p \rfloor$ (zero base-$p^2$ carries).
\end{lemma}

\begin{proof}
For any prime $p \le \sqrt{40n}$, the maximum possible exponent lost in $v_p$ is $O(\log_p(40n))$, contributing $O(\log n)$ to $\log\mathcal{D}_n$; summing over the $\pi(\sqrt{40n}) = O(\sqrt{n}/\log n)$ such primes gives total loss $O(\sqrt{n}) = o(n)$. The second claim is immediate since $p^2 > 40n \ge k$.
\end{proof}

For $x_0 \in \RR$ and $y \in \RR$, define the $1$-periodic (in both $x_0$ and $y$) base-$p$ carry count of the five binomials in~\eqref{eq:aj-binom}:
\begin{equation}
\label{eq:rho}
\begin{aligned}
  \rho(x_0, y) &:= \sum_{r \in \{13, 15, 17\}} \Bigl(\lfloor y + r x_0 \rfloor - \lfloor 13 x_0 \rfloor - \lfloor y + (r - 13) x_0 \rfloor\Bigr) \\
               &\quad + \Bigl(\lfloor 20 x_0 \rfloor - \lfloor y - 3 x_0 \rfloor - \lfloor 23 x_0 - y \rfloor\Bigr)
                      + \Bigl(\lfloor 19 x_0 \rfloor - \lfloor y - 2 x_0 \rfloor - \lfloor 21 x_0 - y \rfloor\Bigr),
\end{aligned}
\end{equation}
and set
\[
  v_{\min}(x_0) := \min_{y \in [0, 1) \setminus \mathcal{B}(x_0)} \rho(x_0, y),
\]
where $\mathcal{B}(x_0) := \{0, \{-2x_0\}, \{-4x_0\}, \{-13x_0\}, \{-15x_0\}, \{-17x_0\}, \{2x_0\}, \{3x_0\}, \{21x_0\}, \{23x_0\}\} \subset [0, 1)$ is the finite set of jump discontinuities (breakpoints) of $\rho(x_0, \cdot)$ on $\RR/\ZZ$.

For a prime $p > \sqrt{40n}$ with $u = n/p$, denote by $\varphi_3(u)$ (for $t_0 = 2, M = 25, \text{order } 3$) and $\varphi_4(u)$ (for $t_0 = 4, M = 27, \text{order } 4$) the baseline termwise carry valuation:
\begin{itemize}
  \item For $u \ge 1$ ($p \le n$): $\varphi_3(u) = \varphi_4(u) = v_{\min}(\{u\})$.
  \item For $u \in (1/M, 1)$ ($n < p < Mn$): $\varphi_s(u)$ is the minimum over $t = j/p \in (2u, 23u)$ of the exact $p$-adic valuation of each summand in $p^s C_n(Q_n, P_{3,n}, \ldots, P_{s,n})$, incorporating both the double-pole band $t \in (3u, 21u)$ and both simple-pole wings $t \in (2u, 3u) \cup (21u, 23u)$, as well as the harmonic jump at $t + t_0 u = 1$.
\end{itemize}

\begin{lemma}[Baseline Carry Divisor]
\label{lem:carry}
For every $n \ge 1$ and every prime $p > \sqrt{40n}$ with $u = n/p > 1/25$ (resp.\ $u > 1/27$):
\begin{enumerate}
  \item[\emph{(a)}] On the entire periodic interval $x_0 \in (0, 1)$, the simple-pole wing carry valuations (plus the $+1$ from $p\,C_n b_j$) are everywhere $\ge v_{\min}(x_0)$ (verified exactly across all $490$ cells of the Farey partition of order $40$).
  \item[\emph{(b)}] Consequently, setting $\mathcal{D}_{s,n}^{\mathrm{car}} := \prod_{\sqrt{40n} < p < M_s n} p^{\varphi_s(n/p)}$ ($M_3 = 25, M_4 = 27$), we have
  \[
    \frac{2\, d_{25n}^3 C_n}{\mathcal{D}_{3,n}^{\mathrm{car}}}(Q_n, P_{3,n}) \in \ZZ^2,
    \qquad
    \frac{6\, d_{27n}^4 C_n}{\mathcal{D}_{4,n}^{\mathrm{car}}}(Q_n, P_{3,n}, P_{4,n}) \in \ZZ^3.
  \]
\end{enumerate}
\end{lemma}

\begin{proof}
Because $p > \sqrt{40n}$, every factorial valuation is simply $v_p(k!) = \lfloor k/p \rfloor$. For $u \ge 1$, every pole satisfies $j + t_0 n \ge 4n \ge 4p$, so $v_p(p^s H_{j+t_0 n}^{(s)}) = 0$, $v_p(p^{s-1} H_{j+t_0 n}^{(s-1)}) = 0$, and $v_p(p L_j) \ge 0$; thus every double-pole term has valuation $\ge \rho(\{u\}, \{j/p\}) \ge v_{\min}(\{u\})$, and by the exact Farey-$40$ check in (a) every simple-pole wing term also has valuation $\ge v_{\min}(\{u\})$. For $u \in (1/M_s, 1)$, $\varphi_s(u)$ is by definition the global minimum over all $t = j/p \in (2u, 23u)$ of all double-pole and simple-pole terms.
\end{proof}

\subsection{The Universal Wilson--Lucas Polynomial-Degree Theorem}
Define the $1$-periodic \emph{Wilson--Lucas defect}
\begin{equation}
\label{eq:Delta}
  \Delta(u) := 2 - 3\lfloor 13u \rfloor + \lfloor 20u \rfloor + \lfloor 19u \rfloor.
\end{equation}
Because $3 \cdot 13 = 20 + 19$, we have $\Delta(u + 1) = \Delta(u)$ identically for all $u \in \RR$ (formally proved in Lean~4, \texttt{wilsonLucasDelta\_periodic}).

Let $\mathcal{S} = \bigcup_{k=1}^{24} I_k \subset (0, 1)$ be the union of the $24$ open intervals listed in Table~\ref{tab:shells24}, which refine along breakpoints $m/k$ ($k \in \{13, 15, 17, 19, 20, 21, 23, 25\}$) into the $38$ open \emph{uniform-carry sub-shells} listed in Appendix~\ref{app:subshells}. Observe why checking the Farey partition $\mathcal{F}_{40}$ of order $40$ is \emph{provably universal}: two breakpoints $\{k_1 x_0\}$ and $\{k_2 x_0\}$ in $\mathcal{B}(x_0)$ (with $k_1, k_2 \in \{-17, -15, -13, -4, -2, \allowbreak 0, 2, 3, 21, 23\}$) can collide modulo $1$ only when $(k_1 - k_2)x_0 \in \ZZ$, and $\max |k_1 - k_2| = 23 - (-17) = 40$. Thus on every open cell of $\mathcal{F}_{40}$, the circular ordering of $\mathcal{B}(x_0)$ and all carry tuples $(c_1,\ldots,c_5)$ are strictly constant! Furthermore, for any prime $p > \max(40, \sqrt{40n})$ with $p \nmid n$, the fractional part $x_0 = \{n/p\}$ has lowest-terms denominator $p > 40$ and hence \emph{never} coincides with any Farey-$40$ endpoint. Exact rational verification over all $88$ open Farey-$40$ cells contained in $\mathcal{S}$ in \texttt{certify\_paper\_bounds.py} (and in Lean~4, \texttt{all\_38\_subshells\_\-defect\_eq\_one}) therefore establishes universally:
\begin{quote}
  \textbf{(Single-Arc Defect Property $\star$):} For every $x_0 \in \mathcal{S} \setminus \mathcal{F}_{40}$ (in particular, for $x_0 = \{n/p\}$ at every prime $p > \max(40, \sqrt{40n})$ with $p \nmid n$ and $\{n/p\} \in \mathcal{S}$):
  \begin{enumerate}
    \item[(i)] $\Delta(x_0) - v_{\min}(x_0) = 1$;
    \item[(ii)] the active minimum set $\mathcal{A}(x_0) := \{y \in \RR/\ZZ \setminus \mathcal{B}(x_0) : \rho(x_0, y) = v_{\min}(x_0)\}$ is a \emph{single contiguous open arc} bounded by two consecutive elements of $\mathcal{B}(x_0)$, on which the five individual binomial carry indicators $(c_1, c_2, c_3, c_4, c_5) \in \{0, 1\}^5$ are constant with $\sum_{i=1}^5 c_i = v_{\min}(x_0)$;
    \item[(iii)] for $u = x_0 \in (1/25, 1)$, $\varphi_3(x_0) = \varphi_4(x_0) = v_{\min}(x_0)$.
  \end{enumerate}
\end{quote}

We can now state and prove the central arithmetic theorem of the paper in half a page.

\begin{theorem}[Universal Wilson--Lucas $\FF_p$-Shell Theorem]
\label{thm:wilson-lucas}
Let $n \ge 1$ and let $p > \max(40, \sqrt{40n})$ be any prime with $p \nmid n$ such that $x_0 := \{n/p\} \in \mathcal{S} = \bigcup_{k=1}^{24} I_k$ (in either the non-periodic regime $u = n/p \in (1/25, 1)$ or the periodic regime $u = n/p \ge 1$). Then $p$ divides every numerator and denominator with \emph{at least one extra power} beyond the carry valuation $v := v_{\min}(x_0)$:
\begin{equation}
\label{eq:extra-p}
  p^{v + 1} \mid C_n Q_n,
  \qquad
  p^{v + 1} \mid 2 p^3 C_n P_{3,n},
  \qquad
  p^{v + 1} \mid 6 p^4 C_n P_{4,n}.
\end{equation}
Consequently, defining the first-order shell divisor
\[
  \mathcal{D}_n^{\mathrm{sh}} := \prod_{\substack{p > \max(40, \sqrt{40n}) \\ \{n/p\} \in \mathcal{S}}} p,
\]
we have for every $n \ge 1$,
\begin{equation}
\label{eq:full-int}
  \frac{2\, d_{25n}^3 C_n}{\mathcal{D}_{3,n}^{\mathrm{car}}\,\mathcal{D}_n^{\mathrm{sh}}}(Q_n, P_{3,n}) \in \ZZ^2,
  \qquad
  \frac{6\, d_{27n}^4 C_n}{\mathcal{D}_{4,n}^{\mathrm{car}}\,\mathcal{D}_n^{\mathrm{sh}}}(Q_n, P_{3,n}, P_{4,n}) \in \ZZ^3.
\end{equation}
\end{theorem}

\begin{proof}
Fix $p > \max(40, \sqrt{40n})$ with $p \nmid n$ and $x_0 = \{n/p\} \in \mathcal{S}$, and let $(c_1, \ldots, c_5) \in \{0,1\}^5$ be the unique active carry tuple on $\mathcal{A}(x_0)$, with $v = \sum_{i=1}^5 c_i = \Delta(x_0) - 1$. Write the base-$p$ expansions $13n = q_{13}p + r_{13}$, $20n = q_{20}p + r_{20}$, and $19n = q_{19}p + r_{19}$ ($0 \le r_k < p$). Notice that $q_{20} + q_{19} - 3q_{13} = \Delta(n/p) - 2 = \Delta(x_0) - 2$.

\medskip
\noindent\textbf{Step 1: Bucket Decomposition and the Canonical Polynomial $\mathcal{R}_m(x) \in \FF_p[x]$.}
Fix $t_0 \in \{2n, 4n\}$ (choosing $t_0 = 2n$ for $Q_n, P_{3,n}$ and $t_0 = 4n$ for $P_{4,n}$) and partition the summation range $2n \le j \le 23n$ into $p$-buckets
\[
  j = m p - t_0 + x, \qquad x \in \{0, 1, \ldots, p - 1\},
\]
so that $j + t_0 = m p + x$ has exact base-$p$ quotient $m = \lfloor (j + t_0)/p \rfloor$ throughout bucket $m$. For each residue $x \in \FF_p$, let $c_i(x) \in \{0, 1\}$ be the actual base-$p$ carry of the $i$-th binomial in~\eqref{eq:aj-binom} at $j = mp - t_0 + x$, so $\sum_{i=1}^5 c_i(x) = \rho(x_0, \{j/p\}) \ge v$.

Define the five rational factors over $\FF_p(x)$ corresponding to the active tuple $(c_1,\ldots,c_5)$:
\begin{itemize}
  \item For each numerator binomial $i \in \{1, 2, 3\}$ (with shift $\alpha_i := -t_0 + (r_i - 13)n \bmod p$, where $(r_1, r_2, r_3) = (13, 15, 17)$):
  \[
    N_i(x) := \begin{cases}
      \displaystyle \prod_{k=1}^{r_{13}} (x + \alpha_i + k), & \text{if } c_i = 0 \quad (\deg = r_{13}), \\[8pt]
      \displaystyle \left[\prod_{k=r_{13}+1}^{p} (x + \alpha_i + k)\right]^{-1}, & \text{if } c_i = 1 \quad (\deg = r_{13} - p).
    \end{cases}
  \]
  \item For each denominator binomial $i \in \{4, 5\}$ (with $(E_4, s_4) = (20, 3)$, $(E_5, s_5) = (19, 2)$, and $\beta_i := -t_0 - s_i n \bmod p$):
  \[
    D_i(x) := \begin{cases}
      \displaystyle \prod_{k=r_{E_i}+1}^{p-1} (k - (x + \beta_i)), & \text{if } c_i = 0 \quad (\deg = p - 1 - r_{E_i}), \\[8pt]
      \displaystyle \left[\prod_{k=0}^{r_{E_i}} (k - (x + \beta_i))\right]^{-1}, & \text{if } c_i = 1 \quad (\deg = (p - 1 - r_{E_i}) - p).
    \end{cases}
  \]
\end{itemize}
Form the product $\mathcal{R}_m(x) := N_1(x) N_2(x) N_3(x) D_4(x) D_5(x) \in \FF_p(x)$. Check its order of vanishing at \emph{every} point $x \in \FF_p$:
\begin{itemize}
  \item If $c_i = 0$, the roots of $N_i$ (resp.\ $D_i$) are precisely the residues $x \in \FF_p$ where $c_i(x) = 1$, so $\ord_x(\text{factor } i) = c_i(x) = c_i(x) - c_i$.
  \item If $c_i = 1$, the poles of $N_i$ (resp.\ $D_i$) are precisely the residues $x \in \FF_p$ where $c_i(x) = 0$, so $\ord_x(\text{factor } i) = -(1 - c_i(x)) = c_i(x) - c_i$.
\end{itemize}
Summing over $i = 1, \ldots, 5$, we obtain at \emph{every} $x \in \FF_p$:
\begin{equation}
\label{eq:ord-nonneg}
  \ord_x(\mathcal{R}_m) = \sum_{i=1}^{5} \bigl(c_i(x) - c_i\bigr) = \rho\bigl(x_0, \{j/p\}\bigr) - v \ge 0.
\end{equation}
Because every pole of $\mathcal{R}_m(x)$ in $\overline{\FF}_p$ lies in $\FF_p$,~\eqref{eq:ord-nonneg} proves that $\mathcal{R}_m(x)$ has \textbf{no poles anywhere} and is therefore a \textbf{polynomial in $\FF_p[x]$}. Its exact degree (verified in Lean~4, theorem \texttt{wilson\_lucas\_degree\_eq\_p\_minus\_2}) is
\begin{equation}
\label{eq:deg-p-2}
  \deg(\mathcal{R}_m)
  = \sum_{i=1}^{3}(r_{13} - c_i p) + \sum_{i=4}^{5}(p - 1 - r_{E_i} - c_i p)
  = \bigl(\Delta(x_0) - v\bigr) p - 2 = p - 2 < p - 1.
\end{equation}
Hence by the elementary finite-field power-sum identity $\sum_{x \in \FF_p} x^d = 0$ ($0 \le d < p - 1$),
\begin{equation}
\label{eq:sum-Rm-zero}
  \sum_{x \in \FF_p} \mathcal{R}_m(x) = 0 \quad \text{in } \FF_p.
\end{equation}

\medskip
\noindent\textbf{Step 2: Exact Proportionality of $C_n a_j$ and $p\,C_n b_j$ to $\mathcal{R}_m(x)$ on Every Bucket $m$.}
Observe that the six critical boundaries $j \in \{2n, 3n, 21n, 23n\}$ and $j + t_0 \equiv 0 \pmod p$ ($t_0 \in \{2n, 4n\}$) correspond to the fractional parts $\{2x_0\}, \{3x_0\}, \{21x_0\}, \{23x_0\}, \{-2x_0\}, \{-4x_0\}$, \textbf{all six of which belong to the breakpoint set $\mathcal{B}(x_0)$} (the latter two being the lower endpoints of $\binom{j+15n}{13n}$ and $\binom{j+17n}{13n}$). By the Single-Arc Property~($\star$), none of these breakpoints lies in the interior of the open active arc $\mathcal{A}(x_0)$. Thus for every bucket $m$, the active indices $\{j = mp - t_0 + x : \{j/p\} \in \mathcal{A}(x_0)\}$ lie \emph{entirely} inside one of four disjoint regimes: (1) outside $[2n, 23n]$ (all terms zero), (2) strictly inside the double-pole band $(3n, 21n)$, (3) strictly inside the left simple-pole wing $(2n, 3n)$, or (4) strictly inside the right simple-pole wing $(21n, 23n)$. Moreover, all ten base-$p$ top-digit quotients $\lfloor(j+cn)/p\rfloor$ are \emph{strictly constant} throughout the active arc of bucket $m$.

By Wilson's theorem, $\frac{(Qp+R)!}{p^Q} \equiv (-1)^Q Q!\, R! \pmod p$ and $R!\,(p - 1 - R)! \equiv (-1)^{R+1} \pmod p$ for all $0 \le Q, R < p$. Applying this to~\eqref{eq:aj-binom} shows that on every double-pole bucket $m$:
\begin{itemize}
  \item Outside $\mathcal{A}(x_0)$, $v_p(C_n a_j) > v$ while $\ord_x(\mathcal{R}_m) > 0$, so both $(C_n a_j/p^v) \bmod p$ and $\mathcal{R}_m(x)$ vanish.
  \item Inside $\mathcal{A}(x_0)$, all top-digit factorials $Q_k!$ and signs $(-1)^{Q_k}$ are constant, while the signs $(-1)^{R_4 + R_5}$ from the two denominator Wilson complements combine to $(-1)^{2j - 5n - (Q_4+Q_5)p} = (-1)^{n + Q_4 + Q_5}$, which is independent of $x$ because $p$ is odd. Collecting the remaining $x$-dependent factors recovers $\mathcal{R}_m(x)$ up to a constant $\kappa_m^{(a)} \in \FF_p$:
  \[
    \frac{C_n a_{mp - t_0 + x}}{p^v} \equiv \kappa_m^{(a)}\, \mathcal{R}_m(x) \pmod p \qquad \forall\, x \in \FF_p.
  \]
  \item For the derivative term $p\,C_n b_j = -(C_n a_j)(p L_j)$ on a double-pole bucket, note that since $40n < p^2$, $p H_N \equiv H_{\lfloor N/p \rfloor} \pmod p$; hence $p L_j \equiv \Lambda_m \pmod p$ is \emph{strictly constant} across the active arc $\mathcal{A}(x_0)$. Outside $\mathcal{A}(x_0)$, $v_p(p\,C_n b_j) \ge v_p(C_n a_j) \ge v + 1$. Thus
  \[
    \frac{p\, C_n b_{mp - t_0 + x}}{p^v} \equiv -\Lambda_m \kappa_m^{(a)}\, \mathcal{R}_m(x) \pmod p \qquad \forall\, x \in \FF_p.
  \]
  \item Finally, on a simple-pole wing bucket (e.g.\ $2n < j < 3n$), the factor $p\,(20n)!\,(3n-1-j)!/(23n-j)!$ in~\eqref{eq:bj-left} has bottom digit $p - 1 - ((j - 3n)\bmod p)$ in $(3n - 1 - j)!$; by Wilson's theorem $(p - 1 - R)! \equiv (-1)^{R+1}/R! \pmod p$, its mod-$p$ dependence on $x$ is \emph{identical} to the fourth binomial $\binom{20n}{j-3n}$ (with the extra factor $p$ compensating the quotient shift $\lfloor(3n-1-j)/p\rfloor = -\lfloor(j-3n)/p\rfloor - 1$), and its alternating sign $(-1)^j$ cancels the remaining denominator Wilson sign $(-1)^{R_5}$. The right wing~\eqref{eq:bj-right} behaves identically. Hence on simple-pole wing buckets as well,
  \[
    \frac{p\, C_n b_{mp - t_0 + x}}{p^v} \equiv \kappa_m^{(b)}\, \mathcal{R}_m(x) \pmod p \qquad \forall\, x \in \FF_p.
  \]
\end{itemize}

\medskip
\noindent\textbf{Step 3: Simultaneous Vanishing of $Q_n$, $P_{3,n}$, and $P_{4,n}$.}
For $j + t_0 = m p + x$ ($0 \le x < p$) and any exponent $s \in \{2, 3, 4\}$:
\[
  p^s H_{mp + x}^{(s)} \equiv H_m^{(s)} \pmod{p^s}
\]
is \textbf{independent of $x$} across the entire bucket $m$ (and for $m = 0$, $p^s H_x^{(s)} \equiv 0 \pmod{p^s}$, contributing an extra factor of $p^s$ automatically). Multiplying the congruences of Step~2 by the bucket constants $H_m^{(s)}$ and summing over $x \in \FF_p$ gives, by~\eqref{eq:sum-Rm-zero},
\begin{align*}
  \sum_{x=0}^{p-1} \frac{C_n a_{mp-t_0+x}}{p^v} &\equiv 0 \pmod p, \\
  \sum_{x=0}^{p-1} \frac{C_n a_{mp-t_0+x}\, p^s H_{mp+x}^{(s)}}{p^v} &\equiv 0 \pmod p, \\
  \sum_{x=0}^{p-1} \frac{p\, C_n b_{mp-t_0+x}\, p^{s-1} H_{mp+x}^{(s-1)}}{p^v} &\equiv 0 \pmod p
\end{align*}
for every bucket $m$ individually! Summing over $m$ proves~\eqref{eq:extra-p}, and combining with Lemma~\ref{lem:carry} (noting Property~($\star$)(iii) that $\varphi_3 = \varphi_4 = v_{\min}$ on $\mathcal{S}$) yields~\eqref{eq:full-int}.
\end{proof}

% =============================================================================
\section{Prime Number Theorem, Digamma Enclosures, and Proofs of Theorems~\ref{thm:main-zeta3} and~\ref{thm:main-joint}}
\label{sec:proof-main}
% =============================================================================

\subsection{Asymptotic Divisor Limits via the Prime Number Theorem}
Recall the Chebyshev form of the Prime Number Theorem: $\sum_{p \le X} \log p = X + o(X)$. Consequently, for any open interval $(a, b) \subset (0, 1)$,
\begin{equation}
\label{eq:pnt-u1}
  \lim_{n \to \infty} \frac{1}{n} \sum_{\substack{p > \sqrt{40n} \\ n/p \in (a, b)}} \log p
  = \frac{1}{a} - \frac{1}{b},
\end{equation}
and for the periodic lifts $n/p \in \bigcup_{k=1}^{\infty} (k + a, k + b)$,
\begin{equation}
\label{eq:pnt-per}
  \lim_{n \to \infty} \frac{1}{n} \sum_{k=1}^{\lfloor \sqrt{n/40} \rfloor} \sum_{n/(k+b) < p < n/(k+a)} \log p
  = \sum_{k=1}^{\infty} \left(\frac{1}{k + a} - \frac{1}{k + b}\right)
  = \psi(b + 1) - \psi(a + 1),
\end{equation}
where $\psi(w) = \Gamma'(w)/\Gamma(w)$ is the digamma function.

\begin{table}[ht]
\centering
\caption{The $24$ merged first-order Wilson--Lucas $\FF_p$-shells $I_k = (a_k, b_k)$ of $F_n^{\star}(t)$, with exact $u < 1$ weight $\frac{1}{a_k} - \frac{1}{b_k}$ and periodic digamma weight $\psi(b_k+1) - \psi(a_k+1)$. (Entries marked $\dagger$ are the two windows of Remark~\ref{rem:2nd-order} where empirical $\ZZ/p^2\ZZ$ vanishing is also observed; Theorems~\ref{thm:main-zeta3} and~\ref{thm:main-joint} use multiplicity $1$ everywhere.)}
\label{tab:shells24}
\footnotesize
\begin{tabular}{@{}clrr@{\hspace{18pt}}clrr@{}}
\toprule
$k$ & $I_k = (a_k, b_k)$ & $\frac{1}{a_k} - \frac{1}{b_k}$ & $\psi(b_k+1)-\psi(a_k+1)$ &
$k$ & $I_k = (a_k, b_k)$ & $\frac{1}{a_k} - \frac{1}{b_k}$ & $\psi(b_k+1)-\psi(a_k+1)$ \\
\midrule
 1 & $(\frac{1}{15}, \frac{2}{27})$ & $1.500000$ & $0.011041$ & 13 & $(\frac{11}{23}, \frac{12}{25})$ & $0.007576$ & $0.001656$ \\
 2 & $(\frac{1}{13}, \frac{2}{25})$ & $0.500000$ & $0.004537$ & 14 & $(\frac{7}{13}, \frac{13}{23})$ & $0.087912$ & $0.023911$ \\
 3 & $(\frac{1}{11}, \frac{2}{21})$ & $0.500000$ & $0.006261$ & 15 & $(\frac{8}{13}, \frac{13}{21})$ & $0.009615$ & $0.003101$ \\
 4 & $(\frac{2}{17}, \frac{3}{25})$ & $0.166667$ & $0.003292$ & 16 & $(\frac{11}{17}, \frac{15}{23})$ & $0.012121$ & $0.004219$ \\
 5 & $(\frac{2}{13}, \frac{4}{25})$ & $0.250000$ & $0.008210$ & 17 & $(\frac{9}{13}, \frac{16}{23})$ & $0.006944$ & $0.002665$ \\
 6 & $(\frac{3}{16}, \frac{4}{21})$ & $0.083333$ & $0.003821$ & 18 & $(\frac{11}{15}, \frac{17}{23})$ & $0.010695$ & $0.004475$ \\
 7 & $(\frac{3}{13}, \frac{6}{25})$ & $0.166667$ & $0.011234$ & 19 & $(\frac{10}{13}, \frac{18}{23})$ & $0.022222$ & $0.010031$ \\
 8 & $(\frac{3}{11}, \frac{7}{25})$ & $0.095238$ & $0.008460$ & 20 & $(\frac{14}{17}, \frac{19}{23})$ & $0.003759$ & $0.001852$ \\
 9 & $(\frac{5}{16}, \frac{8}{25})$ & $0.075000$ & $0.008363$ & 21 & $(\frac{11}{13}, \frac{23}{27})$ & $0.007905$ & $0.004058$ \\
10 & $(\frac{5}{13}, \frac{9}{23})$ & $0.044444$ & $0.006939$ & 22 & $(\frac{14}{15}, \frac{16}{17})$ & $0.008929$ & $0.005265$ \\
11 & $(\frac{6}{13}, \frac{7}{15})$ & $0.023810$ & $0.004951$ & 23$^\dagger$ & $(\frac{16}{17}, \frac{20}{21})$ & $0.012500$ & $0.007475$ \\
12$^\dagger$ & $(\frac{7}{15}, \frac{11}{23})$ & $0.051948$ & $0.011109$ & 24 & $(\frac{20}{21}, \frac{22}{23})$ & $0.004545$ & $0.002749$ \\
\bottomrule
\end{tabular}
\end{table}

\begin{proposition}[Rigorous Certified Enclosures for the Asymptotic Divisors]
\label{prop:enclosures}
Evaluating~\eqref{eq:pnt-u1} in exact rational arithmetic over the $490$-cell Farey partition of order $40$ and enclosing~\eqref{eq:pnt-per} via the strictly alternating Bernoulli asymptotic expansion of $\psi(x + 61)$ through $B_{18}$ (whose remainder has the same sign and smaller magnitude than the first omitted term $-B_{20}/(20 w^{20}) > 0$, see~\cite[Eq.~5.11.2]{DLMF}), we obtain:
\begin{align*}
  D_{3,\mathrm{car}}^{(u<1)} &= \frac{1019365273}{29099070} = 35.0308540101\ldots, \\
  D_{4,\mathrm{car}}^{(u<1)} &= \frac{1193959693}{29099070} = D_{3,\mathrm{car}}^{(u<1)} + 6, \\
  D_{\mathrm{sh}}^{(u<1)} &= \frac{19552742437}{5354228880} = 3.6518316410\ldots, \\
  D_{\mathrm{car}}^{(\mathrm{per})} &\in [1.166093849771775,\, 1.166093849771776], \\
  D_{\mathrm{sh}}^{(\mathrm{per})} &\in [0.159677283334464,\, 0.159677283334465].
\end{align*}
Hence the total first-order asymptotic log-divisors $D_3 := \lim_{n\to\infty}\frac{1}{n}\log(\mathcal{D}_{3,n}^{\mathrm{car}}\mathcal{D}_n^{\mathrm{sh}})$ and $D_4 := \lim_{n\to\infty}\frac{1}{n}\log(\mathcal{D}_{4,n}^{\mathrm{car}}\mathcal{D}_n^{\mathrm{sh}})$ satisfy
\begin{equation}
\label{eq:D3-D4}
  D_3 > 40.00845678419,
  \qquad
  D_4 = D_3 + 6 > 46.00845678419.
\end{equation}
\end{proposition}

\begin{proof}
All $490$ Farey cells are checked in exact rational arithmetic in \texttt{certify\_paper\_bounds.py}; the exact integer identities $D_{4,\mathrm{car}}^{(u<1)} - D_{3,\mathrm{car}}^{(u<1)} = 6$ and $D_{\mathrm{sh}}^{(u<1)} = 19552742437 / 5354228880$ are kernel-verified in Lean~4 (theorems \texttt{exact\_u1\_carry\_joint\_minus\_z3} and \texttt{subshells38\_exact\_u1\_sum}). For the periodic digamma sums, using shift $N = 60$ gives $w = x + 61 \ge 61$ and omitted-term width $\le |B_{20}|/(20 \cdot 61^{20}) < 10^{-34}$ per cell, yielding total enclosure widths $5.16 \times 10^{-32}$ and $2.13 \times 10^{-33}$.
\end{proof}

\subsection{Standard Diophantine Lemma and Completion of Proofs}
We recall the standard elementary lemma converting exponentially convergent rational approximations with bounded-gap non-vanishing into an upper bound on $\mu(\alpha)$.

\begin{lemma}[Bounded-Gap Irrationality Measure Criterion]
\label{lem:diophantine}
Let $\alpha \in \RR$, and let $(\widetilde{p}_n, \widetilde{q}_n) \in \ZZ^2$ ($n \ge 1$) be integer sequences such that for constants $0 < h < g$:
\begin{enumerate}
  \item[\emph{(i)}] $\limsup_{n\to\infty} \frac{1}{n}\log|\widetilde{q}_n| \le h$;
  \item[\emph{(ii)}] $\limsup_{n\to\infty} \frac{1}{n}\log|\widetilde{q}_n \alpha - \widetilde{p}_n| \le -(g - h)$;
  \item[\emph{(iii)}] there exists $N_0$ such that for every $n \ge N_0$, at least one of the two consecutive indices $m \in \{n, n+1\}$ satisfies $|\widetilde{q}_m \alpha - \widetilde{p}_m| \ge e^{-(g - h + o(1)) m} > 0$.
\end{enumerate}
Then $\alpha \notin \QQ$ and $\mu(\alpha) \le \frac{g}{g - h}$.
\end{lemma}

\begin{proof}
Fix any $\varepsilon \in (0, \frac{g - h}{4})$. By~(i)--(iii), there exists $N_1(\varepsilon) \ge N_0$ such that for all $n \ge N_1(\varepsilon)$ we have $|\widetilde{q}_n| \le e^{(h + \varepsilon)n}$ and $|\widetilde{q}_n \alpha - \widetilde{p}_n| \le e^{-(g - h - \varepsilon)n}$, and at least one index $m \in \{n, n+1\}$ satisfies $|\widetilde{q}_m \alpha - \widetilde{p}_m| \ge e^{-(g - h + \varepsilon)m} > 0$.

First, $\alpha \notin \QQ$: if $\alpha = p/q$ with $q \ge 1$, then for the non-vanishing index $m \in \{n, n+1\}$ the positive integer $q|\widetilde{q}_m \alpha - \widetilde{p}_m|$ would tend to $0$ as $n \to \infty$, a contradiction.

Second, let $p/q \in \QQ$ with $\gcd(p, q) = 1$ and $q$ sufficiently large that the minimal integer $n$ with $e^{(g - h - \varepsilon)n} > 2q$ satisfies $n \ge N_1(\varepsilon)$; then $e^{(g - h - \varepsilon)(n - 1)} \le 2q$, so $e^m \le e^2 (2q)^{1/(g - h - \varepsilon)}$ for $m \in \{n, n+1\}$. Let $m \in \{n, n+1\}$ be the index satisfying $|\widetilde{q}_m \alpha - \widetilde{p}_m| \ge e^{-(g - h + \varepsilon)m}$.
\begin{itemize}
  \item \emph{Case 1: $p\,\widetilde{q}_m - q\,\widetilde{p}_m \ne 0$.} Then $1 \le |p\,\widetilde{q}_m - q\,\widetilde{p}_m| = \bigl|\widetilde{q}_m(q\alpha - p) - q(\widetilde{q}_m\alpha - \widetilde{p}_m)\bigr| \le |\widetilde{q}_m|\,|q\alpha - p| + q\,|\widetilde{q}_m\alpha - \widetilde{p}_m|$. Since $q\,|\widetilde{q}_m\alpha - \widetilde{p}_m| \le q\,e^{-(g - h - \varepsilon)n} < \frac{1}{2}$, we get
  \[
    \left|\alpha - \frac{p}{q}\right| > \frac{1}{2q|\widetilde{q}_m|} \ge \frac{1}{2q\, e^{(h + \varepsilon)m}} \gg_\varepsilon q^{-1 - \frac{h+\varepsilon}{g-h-\varepsilon}} = q^{-\frac{g}{g-h-\varepsilon}}.
  \]
  \item \emph{Case 2: $p\,\widetilde{q}_m - q\,\widetilde{p}_m = 0$.} Then $p/q = \widetilde{p}_m/\widetilde{q}_m$, and hence directly by~(i) and~(iii),
  \[
    \left|\alpha - \frac{p}{q}\right| = \frac{|\widetilde{q}_m \alpha - \widetilde{p}_m|}{|\widetilde{q}_m|} \ge \frac{e^{-(g - h + \varepsilon)m}}{e^{(h + \varepsilon)m}} = e^{-(g + 2\varepsilon)m} \gg_\varepsilon q^{-\frac{g + 2\varepsilon}{g - h - \varepsilon}}.
  \]
\end{itemize}
In both cases $|\alpha - p/q| \gg_\varepsilon q^{-\frac{g+2\varepsilon}{g-h-\varepsilon}}$ for all sufficiently large $q$. Letting $\varepsilon \to 0^+$ proves $\mu(\alpha) \le \frac{g}{g - h}$.
\end{proof}

\begin{proof}[Proof of Theorem~\ref{thm:main-zeta3}]
Define the integer sequences
\[
  \widetilde{q}_n := \frac{2\, d_{25n}^3 C_n Q_n}{\mathcal{D}_{3,n}^{\mathrm{car}}\,\mathcal{D}_n^{\mathrm{sh}}} \in \ZZ,
  \qquad
  \widetilde{p}_n := \frac{2\, d_{25n}^3 C_n P_{3,n}}{\mathcal{D}_{3,n}^{\mathrm{car}}\,\mathcal{D}_n^{\mathrm{sh}}} \in \ZZ,
\]
which are integers for all $n \ge 1$ by Theorem~\ref{thm:wilson-lucas}. Since $\frac{1}{n}\log d_{25n}^3 \to 75$ by the Prime Number Theorem, Lemma~\ref{lem:sigma} and Proposition~\ref{prop:enclosures} give
\[
  \limsup_{n\to\infty} \frac{1}{n}\log|\widetilde{q}_n|
  \le \sigma + 75 - D_3
  < 87.053512863 + 75 - 40.008456784
  = 122.045056079 =: h_3.
\]
Meanwhile, $\widetilde{q}_n \zeta(3) - \widetilde{p}_n = \frac{2 d_{25n}^3}{\mathcal{D}_{3,n}^{\mathrm{car}}\mathcal{D}_n^{\mathrm{sh}}} R_{3,n}$, so Proposition~\ref{prop:asymp-nonvanish} (with $A = 1, B = 0$) gives both
\[
  \limsup_{n\to\infty} \frac{1}{n}\log|\widetilde{q}_n\zeta(3) - \widetilde{p}_n|
  \le (75 - D_3) - (g - \sigma)
  = -(g - h_3)
\]
and the bounded-gap lower bound~(iii) of Lemma~\ref{lem:diophantine}. Applying Lemma~\ref{lem:diophantine} with our certified enclosures $g > 150.301021$ and $h_3 < 122.045057$ (Lean~4 \texttt{thm1\_mu\_zeta3\_lt\_53193}) yields
\[
  \mu(\zeta(3)) \le \frac{g}{g - h_3} < \frac{150.301021}{150.301021 - 122.045057} < 5.3193. \qedhere
\]
\end{proof}

\begin{proof}[Proof of Theorem~\ref{thm:main-joint}]
Without loss of generality, scale $(A, B) \in \ZZ^2 \setminus \{(0, 0)\}$ and put $\alpha := A\zeta(3) + B\zeta(4)$. Define the integer sequences
\[
  \widetilde{q}_n^{(4)} := \frac{6\, d_{27n}^4 C_n Q_n}{\mathcal{D}_{4,n}^{\mathrm{car}}\,\mathcal{D}_n^{\mathrm{sh}}} \in \ZZ,
  \qquad
  \widetilde{p}_n^{(4)} := \frac{6\, d_{27n}^4 C_n (A P_{3,n} + B P_{4,n})}{\mathcal{D}_{4,n}^{\mathrm{car}}\,\mathcal{D}_n^{\mathrm{sh}}} \in \ZZ,
\]
which are integers for all $n \ge 1$ by Theorem~\ref{thm:wilson-lucas}. Since $\frac{1}{n}\log d_{27n}^4 \to 108$, Lemma~\ref{lem:sigma} and Proposition~\ref{prop:enclosures} give
\[
  \limsup_{n\to\infty} \frac{1}{n}\log|\widetilde{q}_n^{(4)}|
  \le \sigma + 108 - D_4
  < 87.053512863 + 108 - 46.008456784
  = 149.045056079 =: h_4 < g.
\]
Notice $g - h_4 > 1.255965 > 0$! Because $\widetilde{q}_n^{(4)}\alpha - \widetilde{p}_n^{(4)} = \frac{6 d_{27n}^4}{\mathcal{D}_{4,n}^{\mathrm{car}}\mathcal{D}_n^{\mathrm{sh}}}(A R_{3,n} + B R_{4,n})$, Proposition~\ref{prop:asymp-nonvanish} verifies hypotheses~(ii) and~(iii) of Lemma~\ref{lem:diophantine} for \emph{every} nonzero $(A, B) \in \ZZ^2 \setminus \{(0, 0)\}$. Hence $\alpha = A\zeta(3) + B\zeta(4) \notin \QQ$ (proving $\dim_{\QQ}\langle 1, \zeta(3), \zeta(4)\rangle = 3$), and by Lemma~\ref{lem:diophantine} and Lean~4 \texttt{thm2\_mu\_joint\_lt\_119\_7},
\[
  \mu\bigl(A\zeta(3) + B\zeta(4)\bigr)
  \le \frac{g}{g - h_4}
  < \frac{150.301021}{150.301021 - 149.045057}
  < 119.7. \qedhere
\]
\end{proof}

% =============================================================================
\section{Conjectural Second-Order Refinement and Concluding Remarks}
\label{sec:remarks}
% =============================================================================

\begin{remark}[Empirical Second-Order Windows Giving $\mu(\zeta(3)) < 5.3037$]
\label{rem:2nd-order}
Direct $\ZZ/p^5\ZZ$ testing of $F_n^{\star}(t)$ across $10$ independent primes per cell on both the Farey order-$40$ and order-$55$ partitions (implemented in our Rust certifier \texttt{tools/cmf\_hunter/src/bin/certify\_530.rs} and spot-checked in \texttt{scripts/prove\_fp\_shells\_530.py}) reveals that on the two merged intervals
\[
  I_{12} = \left(\frac{7}{15}, \frac{11}{23}\right),
  \qquad
  I_{23} = \left(\frac{16}{17}, \frac{20}{21}\right),
\]
all three quantities $C_n Q_n$, $2 p^3 C_n P_{3,n}$, and $6 p^4 C_n P_{4,n}$ vanish to \emph{second order} modulo $p^{v + 2}$! Proving this second-order congruence universally (analogous to the $\ZZ/p^2\ZZ$ trace identity used on $(\frac{1}{10}, \frac{2}{19}) \cup (\frac{11}{12}, \frac{15}{16})$ in~\cite{LeibtagKaminer2026}) would add an additional asymptotic divisor weight of
\[
  \left(\frac{15}{7} - \frac{23}{11}\right) + \left(\frac{17}{16} - \frac{21}{20}\right)
  + \sum_{k \in \{12, 23\}} \bigl(\psi(b_k+1) - \psi(a_k+1)\bigr)
  = 0.083032635\ldots,
\]
lowering $h_3$ from $122.045056$ to $121.962023$ and sharpening Theorem~\ref{thm:main-zeta3} from $5.3193$ to
\[
  \mu_{\mathrm{full}}(\zeta(3)) \le 5.303681511 < 5.3037
\]
(and tightening Theorem~\ref{thm:main-joint} to $\mu(A\zeta(3) + B\zeta(4)) < 112.3$). We have deliberately excluded these two unproved second-order windows from Theorems~\ref{thm:main-zeta3} and~\ref{thm:main-joint} so that our main proofs remain $100\%$ self-contained and elementary.
\end{remark}

\begin{remark}[Reproducibility and Software Artifacts]
All verification scripts and Lean~4 proofs are publicly available at \url{https://github.com/aravind-kolumum-raja/zeta3-closed-form}:
\begin{itemize}
  \item \texttt{scripts/certify\_paper\_bounds.py}: standalone pure-Python-stdlib $100$-digit interval-arithmetic certifier for all $38$ sub-shells, saddle roots, exact rational concavity polynomials, and alternating-Bernoulli digamma enclosures.
  \item \texttt{scripts/prove\_fp\_shells\_530.py}: direct mod-$p^k$ verification of the Wilson--Lucas degree identity $\deg(g_m) = p - 2$ and bucket-by-bucket vanishing of $Q_n$, $P_{3,n}$, and $P_{4,n}$ on all $38$ sub-shells in both $u < 1$ and periodic $u = 1 + x$.
  \item \texttt{Zeta3/ElementaryRecordProof.lean}: zero-dependency Lean~4 (v4.36.0) formal verification building via \texttt{lake build} in under one second (reproduced in full in Appendix~\ref{app:lean}).
\end{itemize}
\end{remark}

% =============================================================================
\appendix
\section{All 38 Uniform-Carry Sub-Shells of \texorpdfstring{$F_n^{\star}(t)$}{F\_n*(t)}}
\label{app:subshells}
% =============================================================================

Table~\ref{tab:subshells38} records all $38$ uniform-carry sub-shells partitioning the $24$ merged intervals of Table~\ref{tab:shells24}. On every sub-shell, $\Delta - v_{\min} = 1$, the active set $\mathcal{A}(x_0)$ is a single contiguous open arc on $\RR/\ZZ$, and the unique active carry tuple $(c_1, c_2, c_3, c_4, c_5) \in \{0,1\}^5$ is listed explicitly.

\begin{longtable}{@{}clcccl@{}}
\caption{Exact uniform-carry decomposition of the $24$ shells into $38$ sub-shells.}
\label{tab:subshells38} \\
\toprule
Shell & Sub-shell $(a, b)$ & $v_{\min}$ & $\Delta$ & Active carry tuple $(c_1,\dots,c_5)$ & Arc length (at midpt) \\
\midrule
\endfirsthead
\toprule
Shell & Sub-shell $(a, b)$ & $v_{\min}$ & $\Delta$ & Active carry tuple $(c_1,\dots,c_5)$ & Arc length (at midpt) \\
\midrule
\endhead
\bottomrule
\endfoot
 1 & $(1/15,\ 2/27)$ & $3$ & $4$ & $(1,1,1,0,0)$ & $0.26667$ \\
 2 & $(1/13,\ 2/25)$ & $0$ & $1$ & $(0,0,0,0,0)$ & $0.41231$ \\
 3 & $(1/11,\ 2/21)$ & $0$ & $1$ & $(0,0,0,0,0)$ & $0.13853$ \\
 4 & $(2/17,\ 3/25)$ & $2$ & $3$ & $(0,1,1,0,0)$ & $0.09882$ \\
 5 & $(2/13,\ 3/19)$ & $0$ & $1$ & $(0,0,0,0,0)$ & $0.11741$ \\
 5 & $(3/19,\ 4/25)$ & $1$ & $2$ & $(0,0,0,0,1)$ & $0.13895$ \\
 6 & $(3/16,\ 4/21)$ & $1$ & $2$ & $(1,0,0,0,0)$ & $0.16518$ \\
 7 & $(3/13,\ 4/17)$ & $0$ & $1$ & $(0,0,0,0,0)$ & $0.19457$ \\
 7 & $(4/17,\ 5/21)$ & $0$ & $1$ & $(0,0,0,0,0)$ & $0.21289$ \\
 7 & $(5/21,\ 6/25)$ & $0$ & $1$ & $(0,0,0,0,0)$ & $0.17524$ \\
 8 & $(3/11,\ 7/25)$ & $2$ & $3$ & $(0,1,0,0,1)$ & $0.30182$ \\
 9 & $(5/16,\ 6/19)$ & $0$ & $1$ & $(0,0,0,0,0)$ & $0.22533$ \\
 9 & $(6/19,\ 8/25)$ & $1$ & $2$ & $(0,0,0,0,1)$ & $0.23158$ \\
10 & $(5/13,\ 9/23)$ & $0$ & $1$ & $(0,0,0,0,0)$ & $0.14716$ \\
11 & $(6/13,\ 7/15)$ & $0$ & $1$ & $(0,0,0,0,0)$ & $0.28205$ \\
12 & $(7/15,\ 8/17)$ & $0$ & $1$ & $(0,0,0,0,0)$ & $0.37255$ \\
12 & $(8/17,\ 9/19)$ & $0$ & $1$ & $(0,0,0,0,0)$ & $0.44272$ \\
12 & $(9/19,\ 10/21)$ & $1$ & $2$ & $(0,0,0,0,1)$ & $0.40100$ \\
12 & $(10/21,\ 11/23)$ & $1$ & $2$ & $(0,0,0,0,1)$ & $0.36439$ \\
13 & $(11/23,\ 12/25)$ & $1$ & $2$ & $(0,0,0,0,1)$ & $0.33391$ \\
14 & $(7/13,\ 11/20)$ & $0$ & $1$ & $(0,0,0,0,0)$ & $0.34038$ \\
14 & $(11/20,\ 14/25)$ & $1$ & $2$ & $(0,0,0,1,0)$ & $0.45500$ \\
14 & $(14/25,\ 13/23)$ & $1$ & $2$ & $(0,0,0,1,0)$ & $0.31043$ \\
15 & $(8/13,\ 13/21)$ & $0$ & $1$ & $(0,0,0,0,0)$ & $0.10989$ \\
16 & $(11/17,\ 13/20)$ & $1$ & $2$ & $(0,1,0,0,0)$ & $0.16324$ \\
16 & $(13/20,\ 15/23)$ & $2$ & $3$ & $(0,1,0,1,0)$ & $0.14022$ \\
17 & $(9/13,\ 16/23)$ & $0$ & $1$ & $(0,0,0,0,0)$ & $0.18562$ \\
18 & $(11/15,\ 14/19)$ & $1$ & $2$ & $(0,1,0,0,0)$ & $0.23158$ \\
18 & $(14/19,\ 17/23)$ & $2$ & $3$ & $(0,1,0,0,1)$ & $0.19222$ \\
19 & $(10/13,\ 18/23)$ & $0$ & $1$ & $(0,0,0,0,0)$ & $0.25753$ \\
20 & $(14/17,\ 19/23)$ & $2$ & $3$ & $(1,1,0,0,0)$ & $0.26982$ \\
21 & $(11/13,\ 17/20)$ & $0$ & $1$ & $(0,0,0,0,0)$ & $0.11346$ \\
21 & $(17/20,\ 23/27)$ & $1$ & $2$ & $(0,0,0,1,0)$ & $0.16759$ \\
22 & $(14/15,\ 16/17)$ & $0$ & $1$ & $(0,0,0,0,0)$ & $0.30588$ \\
23 & $(16/17,\ 18/19)$ & $0$ & $1$ & $(0,0,0,0,0)$ & $0.49536$ \\
23 & $(18/19,\ 19/20)$ & $1$ & $2$ & $(0,0,0,0,1)$ & $0.46184$ \\
23 & $(19/20,\ 20/21)$ & $2$ & $3$ & $(0,0,0,1,1)$ & $0.43929$ \\
24 & $(20/21,\ 22/23)$ & $2$ & $3$ & $(0,0,0,1,1)$ & $0.40994$ \\
\end{longtable}

% =============================================================================
\section{Exact Small-\texorpdfstring{$n$}{n} Approximants and Steepest-Descent Verification}
\label{app:numeric}
% =============================================================================

Table~\ref{tab:small-n} summarizes exact rational evaluation of $(C_n Q_n, C_n P_{3,n}, C_n P_{4,n})$ for $n = 1, 2, 3$ to $230$ decimal digits alongside the leading steepest-descent asymptotic prediction of Proposition~\ref{prop:asymp-nonvanish}.

\begin{table}[ht]
\centering
\caption{Exact $230$-digit evaluation of $F_n^{\star}(t)$ for $n = 1, 2, 3$ vs.\ Proposition~\ref{prop:asymp-nonvanish}.}
\label{tab:small-n}
\footnotesize
\setlength{\tabcolsep}{4pt}
\begin{tabular}{@{}ccccccc@{}}
\toprule
$n$ & Digits of $C_n Q_n$ & $|\zeta(3) - P_{3,n}/Q_n|$ & $|\zeta(4) - P_{4,n}/Q_n|$ & Exact $R_{3,n}$ & Steepest-descent $R_{3,n}$ & Ratio \\
\midrule
$1$ & $35$  & $5.189 \times 10^{-66}$  & $6.689 \times 10^{-65}$  & $+1.287 \times 10^{-31}$ & $+1.408 \times 10^{-31}$ & $0.915$ \\
$2$ & $72$  & $1.758 \times 10^{-130}$ & $1.338 \times 10^{-132}$ & $-7.101 \times 10^{-59}$ & $-7.487 \times 10^{-59}$ & $0.949$ \\
$3$ & $110$ & $1.553 \times 10^{-196}$ & $1.944 \times 10^{-195}$ & $+1.797 \times 10^{-87}$ & $+1.816 \times 10^{-87}$ & $0.989$ \\
\bottomrule
\end{tabular}
\end{table}

\clearpage
% =============================================================================
\section{Kernel-Verified Lean 4 Formalization}
\label{app:lean}
% =============================================================================

The following self-contained Lean~4 (v4.36.0) module (\texttt{ElementaryRecordProof.lean}, requiring no external libraries) compiles cleanly with \texttt{lake build} (zero \texttt{sorry}, zero \texttt{native\_decide}) and formally verifies the saddle cubic identity~\eqref{eq:saddle-cubic}, the $1$-periodicity of $\Delta(u)$, the universal Wilson--Lucas degree identity~\eqref{eq:deg-p-2}, all $38$ uniform-carry sub-shell defects, the algebraic strict concavity certificate of Lemma~\ref{lem:concavity}, the exact rational $u < 1$ shell sum, and the final inequalities of Theorems~\ref{thm:main-zeta3} and~\ref{thm:main-joint}.

\begin{lstlisting}
namespace Zeta3.ElementaryRecordProof

-- 1. Exact Primitive Saddle Cubic of F_n*(t)
theorem saddle_cubic_identity (x : Int) :
    x * (x - 2) * (x - 4) * (x + 2) * (x + 3)
      - (x - 13) * (x - 15) * (x - 17) * (x + 23) * (x + 21)
    = 15 * (54 * x ^ 3 - 298 * x ^ 2 - 11879 * x + 106743) := by
  grind

def saddleCubic (x : Int) : Int :=
  54 * x ^ 3 - 298 * x ^ 2 - 11879 * x + 106743

theorem saddle_pole_root_bracket :
    saddleCubic (-16) < 0 /\ saddleCubic (-15) > 0 := by
  decide

-- 2. Periodicity of Delta(u) and Universal Wilson-Lucas Degree Identity
def ratFloor (a : Int) (b : Nat) : Int := a / (b : Int)

def wilsonLucasDelta (a : Int) (b : Nat) : Int :=
  2 - 3 * ratFloor (13 * a) b + ratFloor (20 * a) b + ratFloor (19 * a) b

theorem wilsonLucasDelta_periodic (a : Int) (b : Nat) (hb : b > 0) :
    wilsonLucasDelta (a + b) b = wilsonLucasDelta a b := by
  unfold wilsonLucasDelta ratFloor
  have hbp : (b : Int) /= 0 := by omega
  have h13 : 13 * (a + (b : Int)) = 13 * a + (b : Int) * 13 := by omega
  have h20 : 20 * (a + (b : Int)) = 20 * a + (b : Int) * 20 := by omega
  have h19 : 19 * (a + (b : Int)) = 19 * a + (b : Int) * 19 := by omega
  rw [h13, h20, h19, Int.add_mul_ediv_left _ _ hbp,
      Int.add_mul_ediv_left _ _ hbp, Int.add_mul_ediv_left _ _ hbp]
  omega

theorem wilson_lucas_degree_eq_p_minus_2
    (p n q13 r13 q20 r20 q19 r19 c1 c2 c3 c4 c5 : Int)
    (h13 : 13 * n = q13 * p + r13)
    (h20 : 20 * n = q20 * p + r20)
    (h19 : 19 * n = q19 * p + r19)
    (h_defect : (2 - 3 * q13 + q20 + q19) - (c1 + c2 + c3 + c4 + c5) = 1) :
    (r13 - c1 * p) + (r13 - c2 * p) + (r13 - c3 * p)
      + ((p - 1 - r20) - c4 * p) + ((p - 1 - r19) - c5 * p) = p - 2 := by
  grind

-- 3. All 38 Uniform-Carry Sub-Shells Certified: Delta(mid) - v_min = 1
def subShells38 : List (Nat x Nat x Nat x Nat x Int) :=
  [ (1,15, 2,27, 3),  (1,13, 2,25, 0),  (1,11, 2,21, 0),  (2,17, 3,25, 2),
    (2,13, 3,19, 0),  (3,19, 4,25, 1),  (3,16, 4,21, 1),  (3,13, 4,17, 0),
    (4,17, 5,21, 0),  (5,21, 6,25, 0),  (3,11, 7,25, 2),  (5,16, 6,19, 0),
    (6,19, 8,25, 1),  (5,13, 9,23, 0),  (6,13, 7,15, 0),  (7,15, 8,17, 0),
    (8,17, 9,19, 0),  (9,19,10,21, 1),  (10,21,11,23, 1), (11,23,12,25, 1),
    (7,13,11,20, 0),  (11,20,14,25, 1), (14,25,13,23, 1), (8,13,13,21, 0),
    (11,17,13,20, 1), (13,20,15,23, 2), (9,13,16,23, 0),  (11,15,14,19, 1),
    (14,19,17,23, 2), (10,13,18,23, 0), (14,17,19,23, 2), (11,13,17,20, 0),
    (17,20,23,27, 1), (14,15,16,17, 0), (16,17,18,19, 0), (18,19,19,20, 1),
    (19,20,20,21, 2), (20,21,22,23, 2) ]

def checkSubShell38 (entry : Nat x Nat x Nat x Nat x Int) : Bool :=
  let (an, ad, bn, bd, vmin) := entry
  an * bd < bn * ad &&
  (wilsonLucasDelta ((an * bd + bn * ad : Nat) : Int) (2 * ad * bd) - vmin == 1)

theorem all_38_subshells_defect_eq_one :
    subShells38.length = 38 /\ subShells38.all checkSubShell38 = true := by
  decide

-- 4. Algebraic Strict Concavity of Vertical Contour psi''(y) < 0 at c = 43/4
def polyMul (p q : List Int) : List Int :=
  let n := p.length + q.length - 1
  List.ofFn (fun (idx : Fin n) =>
    let k := idx.val
    (List.range p.length).foldl (fun acc i =>
      if i <= k /\ k - i < q.length then
        acc + (p.getD i 0) * (q.getD (k - i) 0)
      else acc) 0)

def contourTermsAt43_4 : List (Int x Int) :=
  [ (-1, 43), (1, -9), (-1, 35), (1, -17), (-1, 27), (1, -25),
    (1, 135), (-1, 51), (1, 127), (-1, 55) ]

def contourConcavityCoeffs : List Int :=
  let n := contourTermsAt43_4.length
  (List.range n).foldl (fun acc i =>
    let (eps, ai) := contourTermsAt43_4.getD i (0, 0)
    let prod := (List.range n).foldl (fun pr j =>
      if j == i then pr
      else polyMul pr [(contourTermsAt43_4.getD j (0,0)).2 ^ 2, 1]) [1]
    List.zipWith (fun x y => x + eps * ai * y) acc prod) (List.replicate n 0)

theorem contour_strictly_concave_at_43_4 :
    contourConcavityCoeffs.getD 9 1 = 0 /\
    (contourConcavityCoeffs.take 9).all (. < 0) = true := by
  decide

-- 5. Exact Rational u < 1 Sums and Final Bounds (Zero Axioms!)
def L23 : Nat := 5354228880

theorem exact_u1_carry_joint_minus_z3 :
    (1193959693 : Int) - 1019365273 = 6 * 29099070 := by decide

theorem subshells38_exact_u1_sum :
    subShells38.all (fun e => L23 % e.1 == 0 && L23 % e.2.2.1 == 0) = true /\
    subShells38.foldl (fun acc e => acc + (L23 / e.1) * e.2.1) 0 =
    subShells38.foldl (fun acc e => acc + (L23 / e.2.2.1) * e.2.2.2.1) 0
      + 19552742437 := by decide

theorem thm1_mu_zeta3_lt_53193 :
    150301021 * 10000 < 53193 * (150301021 - 122045057) := by decide

theorem thm2_mu_joint_lt_119_7 :
    149045057 < 150301021 /\
    150301021 * 10 < 1197 * (150301021 - 149045057) := by decide

end Zeta3.ElementaryRecordProof
\end{lstlisting}
\noindent\emph{(Note: ASCII symbols \texttt{/\textbackslash}, \texttt{/=}, \texttt{<=}, \texttt{x} in the listing above correspond verbatim to Unicode \texttt{$\wedge$}, \texttt{$\ne$}, \texttt{$\le$}, \texttt{$\times$} in the compiled source file \texttt{Zeta3/ElementaryRecordProof.lean}.)}

\bibliographystyle{amsplain}
\bibliography{references}

\end{document}
